\documentclass[10pt,letterpaper]{article}

\usepackage[letterpaper,margin=1in]{geometry}
\usepackage{lmodern}
\usepackage{amsmath,amssymb,amsthm,mathtools}
\usepackage{xcolor}
\usepackage{hyperref}
\usepackage{microtype}

\hypersetup{
  colorlinks=true,
  linkcolor=blue!55!black,
  citecolor=green!50!black,
  urlcolor=blue!55!black,
  pdftitle={A Proof of the Ma--Yang--Li Signature Conjecture via Weighted Signed Graphs},
  pdfauthor={Prince Anand}
}

% Match section and subsection heading sizes to 12pt bold / 10pt bold
\makeatletter
\renewcommand\section{\@startsection{section}{1}{\z@}%
  {-3.2ex \@plus -1ex \@minus -.2ex}%
  {1.8ex \@plus .2ex}%
  {\normalfont\large\bfseries}}
\renewcommand\subsection{\@startsection{subsection}{2}{\z@}%
  {-2.6ex \@plus -0.8ex \@minus -.2ex}%
  {1.2ex \@plus .2ex}%
  {\normalfont\normalsize\bfseries}}
% Indented 9pt abstract with bold run-in "Abstract."
\renewenvironment{abstract}{%
  \begin{list}{}{%
    \setlength{\leftmargin}{0.45in}%
    \setlength{\rightmargin}{0.45in}%
    \setlength{\topsep}{0pt}%
    \setlength{\partopsep}{0pt}%
  }%
  \item\relax\small\textbf{Abstract.}\ %
}{%
  \end{list}\vspace{0.8em}%
}
\makeatother

% Theorem environments: bold upright label, upright parenthetical note, italic body
\newtheoremstyle{zetaplain}%
  {6pt plus 2pt minus 1pt}% Space above
  {6pt plus 2pt minus 1pt}% Space below
  {\itshape}% Body font
  {}% Indent amount
  {\bfseries}% Theorem head font
  {.}% Punctuation after theorem head
  {0.5em}% Space after theorem head
  {\thmname{#1}\thmnumber{ #2}\thmnote{ {\normalfont(#3)}}}% Theorem head spec

\newtheoremstyle{zetadef}%
  {6pt plus 2pt minus 1pt}% Space above
  {6pt plus 2pt minus 1pt}% Space below
  {\normalfont}% Body font
  {}% Indent amount
  {\bfseries}% Theorem head font
  {.}% Punctuation after theorem head
  {0.5em}% Space after theorem head
  {\thmname{#1}\thmnumber{ #2}\thmnote{ {\normalfont(#3)}}}% Theorem head spec

\theoremstyle{zetaplain}
\newtheorem{theorem}{Theorem}[section]
\newtheorem{lemma}[theorem]{Lemma}
\newtheorem{proposition}[theorem]{Proposition}
\newtheorem{corollary}[theorem]{Corollary}

\theoremstyle{zetadef}
\newtheorem{definition}[theorem]{Definition}

\theoremstyle{remark}
\newtheorem*{remark}{Remark}
\newtheorem*{remarks}{Remarks}
\newtheorem*{example}{Example}

\numberwithin{equation}{section}

\DeclareMathOperator{\sgn}{sgn}
\DeclareMathOperator{\im}{im}
\DeclareMathOperator{\tr}{tr}
\DeclareMathOperator{\rank}{rank}
\DeclareMathOperator{\adj}{adj}
\DeclareMathOperator{\diag}{diag}

\newcommand{\vecC}[1]{\overrightarrow{#1}}
\newcommand{\revC}[1]{\overleftarrow{#1}}

\title{\LARGE A Proof of the Ma--Yang--Li Signature Conjecture\\via Weighted Signed Graphs}
\author{\large Prince Anand}
\date{\normalsize November 2026}

\begin{document}
\maketitle

\begin{abstract}
We prove that every finite simple graph $G$ satisfies $-c_3(G) \le s(G) \le c_5(G)$, where $s(G)$ is the adjacency signature and $c_3(G), c_5(G)$ count cycles of lengths $3$ and $1$ modulo $4$. This follows from a general theorem for real symmetric zero-diagonal matrices: if every odd cycle $C$ of length $2k+1$ in the underlying signed graph $(G,\sigma)$ satisfies $(-1)^k \prod_{e\in C}\sigma(e) = -1$, then $s(A) \le 0$. We classify $3$-connected non-bipartite such graphs as $(K_4,+1)$ or $(\mathrm{Petersen},-1)$ up to switching, exclude bipartite SPQR leaves, and eliminate weighted $K_4$- and Petersen-leaves via Schur complements.
\end{abstract}

% ============================================================================
\section{Introduction}\label{sec:intro}
% ============================================================================

Let $A$ be a real symmetric $n \times n$ matrix. Its \emph{inertia} is the triple $\bigl(n^+(A), n^0(A), n^-(A)\bigr)$, where $n^+(A)$, $n^0(A)$, and $n^-(A)$ denote the numbers of positive, zero, and negative eigenvalues of $A$, counted with multiplicity. The \emph{signature} of $A$ is the difference
\[
  s(A) = n^+(A) - n^-(A).
\]
When $G$ is a finite simple graph with adjacency matrix $A(G)$, we write $n^+(G) = n^+(A(G))$, $n^-(G) = n^-(A(G))$, and $s(G) = s(A(G))$.

If $G$ is bipartite, its adjacency spectrum is symmetric about the origin, so $s(G) = 0$. For a general graph, the signature $s(G)$ measures the spectral asymmetry created by odd cycles. The simplest non-bipartite examples already reveal a striking modulo-$4$ pattern. The adjacency eigenvalues of an odd cycle $C_\ell$ of length $\ell$ are $2\cos(2\pi j / \ell)$ for $0 \le j < \ell$, and counting how many lie in $(0, 2]$ versus $[-2, 0)$ yields
\[
  s(C_\ell) =
  \begin{cases}
    -1 & \text{if } \ell \equiv 3 \pmod 4, \\
    +1 & \text{if } \ell \equiv 1 \pmod 4.
  \end{cases}
\]
Let $c_3(G)$ and $c_5(G)$ denote the numbers of cycles in $G$ of lengths congruent to $3$ and $1$ modulo $4$, respectively. In 2013, Ma, Yang, and Li~\cite{MYL13} proved that the two-sided bound
\begin{equation}\label{eq:MYL}
  -c_3(G) \le s(G) \le c_5(G)
\end{equation}
holds for trees, unicyclic graphs, and bicyclic graphs, and conjectured that~\eqref{eq:MYL} holds for every finite simple graph $G$. Over the ensuing decade, the conjecture was confirmed for several restricted classes—including line graphs and power trees by Wang and Fan~\cite{WF14}, and chordal graphs and cographs by He, Qian, Shan, and Wu~\cite{HQSW21}—while the general problem remained open (see, for example, Akbari, Elphick, Kumar, Pragada, and Tang~\cite[Conjecture~2.5]{AEKPT26}).

\medskip
\noindent\textbf{Why weighted signed graphs?}
A direct attack on~\eqref{eq:MYL} for unweighted $0$--$1$ adjacency matrices stalls almost immediately. To prove $s(G) \le c_5(G)$ by induction on $c_5(G)$, one may delete a vertex on a $(1 \bmod 4)$-cycle: Cauchy's interlacing theorem changes the signature by at most $1$, reducing the entire upper bound to the base case $c_5(G) = 0$, where one must prove $s(G) \le 0$. Similarly, the lower bound $-c_3(G) \le s(G)$ reduces to proving $s(-A(G)) \le 0$ whenever $c_3(G) = 0$. However, even when $c_5(G) = 0$, eliminating a degree-$2$ vertex or a small triconnected leaf by a Schur complement alters edge weights and can flip the signs of edges. Unweighted graphs are not closed under these matrix reductions, whereas real symmetric zero-diagonal matrices are.

This observation dictates the right level of generality. Rather than restricting to $\{0, 1\}$-matrices, we work with arbitrary real symmetric zero-diagonal matrices and track only the \emph{signs} of their nonzero entries.

\begin{definition}[signed cycle type and NOC matrices]\label{def:NOC}
Let $A = (a_{ij})$ be a real symmetric $n \times n$ matrix with $a_{ii} = 0$ for all $1 \le i \le n$. Its \emph{underlying signed graph} is the pair $(G, \sigma)$, where $G = (V, E)$ is the simple graph on $V = \{1, \dots, n\}$ with $ij \in E$ if and only if $a_{ij} \ne 0$, and $\sigma : E \to \{+1, -1\}$ is given by $\sigma(ij) = \sgn(a_{ij})$.

For every odd cycle $C$ of $G$ of length $2k+1$, define its \emph{signed cycle type} by
\[
  \epsilon(C) = (-1)^k \prod_{e \in C} \sigma(e) \in \{+1, -1\}.
\]
We say that $(G, \sigma)$ (and $A$) is \emph{negative-odd-cycle}, abbreviated \emph{NOC}, if $\epsilon(C) = -1$ for every odd cycle $C$ of $G$.
\end{definition}

\noindent\textbf{The mod-$4$ weight interpretation.}
It is convenient to translate $\epsilon(C)$ into an additive weight in $\mathbb{Z}/4\mathbb{Z}$. Assign each edge $e \in E$ the weight
\[
  w_\sigma(e) =
  \begin{cases}
    1 \pmod 4 & \text{if } \sigma(e) = +1, \\
    3 \pmod 4 & \text{if } \sigma(e) = -1,
  \end{cases}
\]
and extend $w_\sigma$ additively to any path or cycle $H$ by $w_\sigma(H) = \sum_{e \in E(H)} w_\sigma(e) \pmod 4$. If an odd cycle $C$ of length $2k+1$ has $r$ negative edges, then
\[
  w_\sigma(C) \equiv (2k + 1) + 2r \pmod 4.
\]
Hence $\epsilon(C) = -1$ if and only if $w_\sigma(C) \equiv 3 \pmod 4$, and $\epsilon(C) = +1$ if and only if $w_\sigma(C) \equiv 1 \pmod 4$. Notice what this unifies:
\begin{itemize}
  \item For an all-positive matrix ($\sigma \equiv +1$), the NOC condition means precisely that $G$ has no cycle of length $1 \pmod 4$, i.e., $c_5(G) = 0$.
  \item For an all-negative matrix ($\sigma \equiv -1$), each edge has weight $3 \equiv -1 \pmod 4$, so an odd cycle of length $2k+1$ has weight $-(2k+1) \pmod 4$; thus the NOC condition means precisely that $G$ has no cycle of length $3 \pmod 4$, i.e., $c_3(G) = 0$.
\end{itemize}
In particular, both halves of the Ma--Yang--Li conjecture become special cases of a single statement.

\begin{theorem}[main theorem]\label{thm:main}
Let $A$ be a real symmetric $n \times n$ matrix with zero diagonal whose underlying signed graph $(G, \sigma)$ is NOC. Then $s(A) \le 0$.
\end{theorem}

\begin{corollary}[Ma--Yang--Li conjecture~\cite{MYL13}]\label{cor:MYL}
For every finite simple graph $G$,
\[
  -c_3(G) \le s(G) \le c_5(G).
\]
\end{corollary}

\begin{example}
Two $3$-connected signed graphs will play a distinguished role throughout the paper. First, consider $(K_4, +1)$, the complete graph on $4$ vertices with all edges positive: every odd cycle is a triangle of weight $1+1+1 = 3 \pmod 4$, so $(K_4, +1)$ is NOC; its unweighted adjacency spectrum is $\{3, -1, -1, -1\}$, giving $s(K_4) = 1 - 3 = -2 \le 0$. Second, consider $(\mathrm{Petersen}, -1)$, the Petersen graph with all edges negative: its odd cycles have lengths $5$ and $9$, and with all weights equal to $3 \pmod 4$ we have $5 \cdot 3 \equiv 3 \pmod 4$ and $9 \cdot 3 \equiv 3 \pmod 4$, so $(\mathrm{Petersen}, -1)$ is NOC; its unweighted negative adjacency spectrum is $\{-3^{(1)}, -1^{(5)}, 2^{(4)}\}$, giving $s(-A(\mathrm{Petersen})) = 4 - 6 = -2 \le 0$. Remarkably, these are the \emph{only} $3$-connected non-bipartite NOC signed graphs up to switching, and for \emph{arbitrary} real edge weights conforming to these signs, the signature never becomes positive.
\end{example}

\medskip
\noindent\textbf{Overview of the proof.}
We argue by examining a counterexample $A$ to Theorem~\ref{thm:main} of minimum order $n = |V(G)|$. The proof consists of four interlocking stages, each corresponding to one section of the paper:
\begin{enumerate}
  \item \emph{Local algebraic reductions (Section~\ref{sec:prelim}).} Using Sylvester's law of inertia and Weyl's monotonicity theorem, we prove that any minimum-order counterexample is $2$-connected (Lemma~\ref{lem:cutvertex}) and has no vertex of degree $2$ (Lemma~\ref{lem:deg2}), hence $\delta(G) \ge 3$.
  \item \emph{Classification of $3$-connected non-bipartite NOC graphs (Section~\ref{sec:3conn}).} Adapting the maximal even-cycle framework of Gao, Li, Ma, and Xie~\cite{GLMX24} and Bai, Grzesik, Li, and Prorok~\cite{BGLP26} from unweighted cycle lengths to mod-$4$ signed weights, we prove that every $3$-connected non-bipartite NOC signed graph is switching-equivalent to $(K_4, +1)$ or $(\mathrm{Petersen}, -1)$ (Theorem~\ref{thm:3conn}).
  \item \emph{Elimination of bipartite SPQR leaves (Section~\ref{sec:bipartite}).} When $G$ is $2$-connected with $\delta(G) \ge 3$ but not $3$-connected, its SPQR tree has at least two leaves, each an $R$-node whose skeleton is a $3$-connected simple graph $T = H + xy$ with one virtual edge $xy$. We prove that if $H$ is bipartite, then for every signing $\sigma$ of $H$, there exist two $(x, y)$-paths in $H$ whose mod-$4$ weights differ by $2 \pmod 4$ (Proposition~\ref{prop:bipartite}). This forces every leaf $R$-node of a non-bipartite NOC graph to be non-bipartite, and hence switching-equivalent to $(K_4 - xy, +1)$ or $(\mathrm{Petersen} - xy, -1)$.
  \item \emph{Schur-complement reductions for weighted $K_4$ and Petersen leaves (Section~\ref{sec:schur}).} We prove that any weighted $(K_4, +1)$ or $(\mathrm{Petersen}, -1)$ has $s(A) \le 0$, and that any weighted $(K_4 - xy, +1)$ or $(\mathrm{Petersen} - xy, -1)$ leaf can be eliminated by a Schur complement (and Weyl's monotonicity theorem) to produce a strictly smaller NOC matrix without decreasing the signature (Lemmas~\ref{lem:K4} and~\ref{lem:Petersen}). For the Petersen leaf, the $8$-vertex interior splits into an independent $4$-set $U$ and a $2K_2$ subgraph $W$; this yields a $4 \times 4$ block decomposition whose cofactors have uniform sign patterns when nonsingular, and a rank-$7$ kernel reduction of signature $-1$ when singular.
\end{enumerate}
In Section~\ref{sec:proof}, we combine these four components via the SPQR-tree decomposition to complete the proofs of Theorem~\ref{thm:main} and Corollary~\ref{cor:MYL}.

% ============================================================================
\section{Preliminaries and Algebraic Reductions}\label{sec:prelim}
% ============================================================================

Throughout the paper, all graphs are finite, undirected, and simple unless stated otherwise. For a graph $G$ and $S \subseteq V(G)$, we write $G[S]$ for the subgraph induced by $S$, $G - S = G[V(G) \setminus S]$, and $N_G(v)$ (or $N_H(v) = N_G(v) \cap V(H)$) for the neighborhood of $v$. For standard matrix inequalities—Sylvester's law of inertia, Cauchy's interlacing theorem, and Weyl's monotonicity theorem—we refer to Horn and Johnson~\cite{HJ13}.

\subsection{Switching and congruence invariance}

Let $A$ be a real symmetric zero-diagonal matrix with underlying signed graph $(G, \sigma)$, and let $D = \diag(d_1, \dots, d_n)$ be an invertible real diagonal matrix. By Sylvester's law of inertia, $A$ and $DAD$ have the same inertia and therefore the same signature: $s(DAD) = s(A)$. The underlying graph of $DAD$ is still $G$, while its signing $\sigma'$ satisfies
\[
  \sigma'(uv) = \sgn(d_u)\sgn(d_v)\sigma(uv) \qquad \text{for every } uv \in E(G).
\]
When $d_i \in \{+1, -1\}$, replacing $\sigma$ by $\sigma'$ is called \emph{switching} at the vertex set $\{i \in V(G) : d_i = -1\}$, and two signings related by switching are \emph{switching-equivalent}. Since every cycle $C$ visits each vertex either zero or two times, $\prod_{e \in C} \sigma'(e) = \prod_{e \in C} \sigma(e)$. Consequently, switching preserves the signed cycle type $\epsilon(C)$ and the mod-$4$ weight $w_\sigma(C)$ of every cycle $C$—and in particular preserves the NOC property—while diagonal congruence $A \mapsto DAD$ preserves both the NOC property and the signature $s(A)$.

Notice also a useful consequence for bipartite graphs: if $G$ is bipartite with bipartition $(X, Y)$, switching at $X$ negates every edge, sending $A$ to $-A$. Therefore $s(A) = s(-A) = -s(A)$, which forces $s(A) = 0$.

\subsection{Cut-vertex and degree-\texorpdfstring{$2$}{2} reductions}

We now show that a minimum-order counterexample to Theorem~\ref{thm:main} cannot have a cut-vertex or a vertex of degree $2$. Both reductions use congruence transformations to decouple a small block of vertices from the rest of the graph.

\medskip
\noindent\textbf{Cut-vertices.}
Suppose $v$ is a cut-vertex separating $G_1 - v$ from $G_2 - v$. If the coupling vector from $v$ into $G_1 - v$ does not lie in the column space of $A(G_1 - v)$, a single congruence operation zeroes out the coupling from $v$ into $G_2 - v$, splitting $A$ into a direct sum $A(G_1) \oplus A(G_2 - v)$. Otherwise, both couplings lie in the respective column spaces, and a Schur complement reduces $v$ to a single scalar diagonal entry equal to the sum of the two separate Schur scalars from $G_1$ and $G_2$.

\begin{lemma}[cut-vertex reduction]\label{lem:cutvertex}
Let $A$ be a real symmetric zero-diagonal matrix with underlying NOC signed graph $(G, \sigma)$. Suppose $v \in V(G)$ is a cut-vertex decomposing $G = G_1 \cup G_2$ with $V(G_1) \cap V(G_2) = \{v\}$ and $|V(G_i)| < |V(G)|$ for $i = 1, 2$. If Theorem~\ref{thm:main} holds for all NOC matrices of order strictly less than $|V(G)|$, then $s(A) \le 0$.
\end{lemma}

\begin{proof}
Ordering the vertices as $\{v\}$, $V(G_1) \setminus \{v\}$, and $V(G_2) \setminus \{v\}$, write
\[
  A = \begin{pmatrix}
    0 & r_1^T & r_2^T \\
    r_1 & A_1 & 0 \\
    r_2 & 0 & A_2
  \end{pmatrix},
\]
where $A_i$ is the principal submatrix on $V(G_i) \setminus \{v\}$ (nonempty since $|V(G_i)| < |V(G)|$ and $V(G_1) \cap V(G_2) = \{v\}$). The graphs $G_1$, $G_2$, $G_1 - v$, and $G_2 - v$ are induced subgraphs of $G$, hence all NOC of order strictly less than $|V(G)|$. By hypothesis, $s(A(G_i)) \le 0$ and $s(A_i) \le 0$ for $i = 1, 2$.

\emph{Case 1: $r_1 \notin \im(A_1)$ or $r_2 \notin \im(A_2)$.} By symmetry, assume $r_1 \notin \im(A_1)$. Since $A_1$ is symmetric, $\im(A_1)^\perp = \ker(A_1)$, so there exists $u \in \ker(A_1)$ with $u^T r_1 = 1$. Set
\[
  P = \begin{pmatrix}
    1 & 0 & 0 \\
    0 & I & -u r_2^T \\
    0 & 0 & I
  \end{pmatrix}.
\]
Then $P$ is invertible and, using $A_1 u = 0$ and $r_1^T u = 1$, block multiplication yields
\[
  P^T A P = \begin{pmatrix}
    0 & r_1^T & 0 \\
    r_1 & A_1 & 0 \\
    0 & 0 & A_2
  \end{pmatrix} = A(G_1) \oplus A_2.
\]
By Sylvester's law of inertia, $s(A) = s(A(G_1)) + s(A_2) \le 0 + 0 = 0$.

\emph{Case 2: $r_1 \in \im(A_1)$ and $r_2 \in \im(A_2)$.} Choose vectors $x_1, x_2$ with $A_1 x_1 = r_1$ and $A_2 x_2 = r_2$. Taking
\[
  P = \begin{pmatrix}
    1 & 0 & 0 \\
    -x_1 & I & 0 \\
    -x_2 & 0 & I
  \end{pmatrix},
\]
we obtain $P^T A P = (-r_1^T x_1 - r_2^T x_2) \oplus A_1 \oplus A_2$. Thus
\[
  s(A) = s(A_1) + s(A_2) + s(-r_1^T x_1 - r_2^T x_2),
\]
and analogously $s(A(G_i)) = s(A_i) + s(-r_i^T x_i) \le 0$ for $i = 1, 2$. If $s(A_1) \le -1$ or $s(A_2) \le -1$, then $s(-r_1^T x_1 - r_2^T x_2) \le 1$ immediately gives $s(A) \le -1 + 0 + 1 = 0$. Otherwise $s(A_1) = s(A_2) = 0$, which forces $s(-r_i^T x_i) = s(A(G_i)) \le 0$, hence $-r_i^T x_i \le 0$ for $i = 1, 2$. Consequently $-r_1^T x_1 - r_2^T x_2 \le 0$ and $s(A) \le 0$.
\end{proof}

\noindent\textbf{Degree-$2$ vertices.}
When $v$ has degree $2$ with neighbors $x, y$, scaling $x$ and $y$ makes $a_{vx} = a_{vy} = 1$. Changing basis on $\{x, y\}$ to $e_x + e_y$ and $e_x - e_y$ isolates a $2 \times 2$ hyperbolic block on $\{v, e_x + e_y\}$ of signature $0$ whose Schur complement leaves the rest of the matrix unchanged except for merging $x$ and $y$ into a single vertex $z = e_x - e_y$ with diagonal entry $-2a_{xy}$. The NOC condition on the triangle $vxy$ (if $a_{xy} \ne 0$) forces $a_{xy} > 0$, so this diagonal entry is non-positive and can be raised to $0$ without decreasing the signature by Weyl's monotonicity theorem. Crucially, the resulting edge weights $a_{xw} - a_{yw}$ may take arbitrary real values, which is precisely why Theorem~\ref{thm:main} is formulated for general weighted signed graphs.

\begin{lemma}[degree-$2$ vertex reduction]\label{lem:deg2}
Let $A$ be a real symmetric zero-diagonal matrix of order $n \ge 3$ with underlying NOC signed graph $(G, \sigma)$. Suppose $v \in V(G)$ has degree $2$ in $G$. If Theorem~\ref{thm:main} holds for all NOC matrices of order $n - 2$, then $s(A) \le 0$.
\end{lemma}

\begin{proof}
Let $x, y$ be the two neighbors of $v$ in $G$, and set $U = V(G) \setminus \{v, x, y\}$. Applying an invertible diagonal congruence scaling $x$ and $y$ (which preserves $s(A)$ and the NOC property), we may assume $a_{vx} = a_{vy} = 1$. Ordering the vertices as $v, x, y, U$, we have
\[
  A = \begin{pmatrix}
    0 & 1 & 1 & 0 \\
    1 & 0 & a_{xy} & r_x^T \\
    1 & a_{xy} & 0 & r_y^T \\
    0 & r_x & r_y & A_U
  \end{pmatrix}.
\]
Change basis on $\{v, x, y\}$ via $e'_v = e_v$, $e_+ = e_x + e_y$, $e_- = e_x - e_y$, keeping $U$ fixed. With respect to $(e'_v, e_+, e_-, U)$, $A$ is congruent to
\[
  \tilde{A} = \begin{pmatrix}
    0 & 2 & 0 & 0 \\
    2 & 2a_{xy} & 0 & r_x^T + r_y^T \\
    0 & 0 & -2a_{xy} & r_x^T - r_y^T \\
    0 & r_x + r_y & r_x - r_y & A_U
  \end{pmatrix}.
\]
The leading $2 \times 2$ block $B = \begin{pmatrix} 0 & 2 \\ 2 & 2a_{xy} \end{pmatrix}$ has $\det(B) = -4 < 0$, so $B$ is nonsingular with $s(B) = 0$ and inverse $B^{-1} = \frac{1}{4}\begin{pmatrix} -2a_{xy} & 2 \\ 2 & 0 \end{pmatrix}$. Because the $(2, 2)$ entry of $B^{-1}$ is $0$ and the first row of the off-diagonal block $R = \begin{pmatrix} 0 & 0 \\ 0 & r_x^T + r_y^T \end{pmatrix}$ is identically $0$, the product $R^T B^{-1} R$ vanishes identically. Hence the Schur complement of $B$ in $\tilde{A}$ is simply the lower-right $(n - 2) \times (n - 2)$ block:
\[
  A' = \begin{pmatrix}
    -2a_{xy} & r_x^T - r_y^T \\
    r_x - r_y & A_U
  \end{pmatrix},
  \qquad s(A) = s(B) + s(A') = s(A').
\]

We claim that $a_{xy} \ge 0$. Indeed, if $a_{xy} \ne 0$, then $(v, x, y)$ is a triangle in $G$ with edge signs $(+1, +1, \sgn(a_{xy}))$. Because $(G, \sigma)$ is NOC, $-1 = \epsilon(vxy) = (-1)^1 (+1)(+1)\sgn(a_{xy})$, forcing $a_{xy} > 0$.

If $n = 3$, then $U = \emptyset$ and $A' = (-2a_{xy})$ is $1 \times 1$ with $-2a_{xy} \le 0$, so $s(A) = s(A') \le 0$. Assume now that $n \ge 4$, and let $A'_0$ be the $(n - 2) \times (n - 2)$ zero-diagonal matrix obtained from $A'$ by replacing its $(1, 1)$ entry $-2a_{xy}$ with $0$. Since $-2a_{xy} \le 0$, we have $A' = A'_0 - 2a_{xy} e_1 e_1^T \preceq A'_0$. By Weyl's monotonicity theorem, every eigenvalue of $A'$ is at most the corresponding eigenvalue of $A'_0$, so $n^+(A') \le n^+(A'_0)$ and $n^-(A') \ge n^-(A'_0)$, giving
\[
  s(A) = s(A') \le s(A'_0).
\]

It remains to verify that the underlying signed graph $(G', \sigma')$ of $A'_0$ (on vertex set $\{z\} \cup U$, where $z$ corresponds to the first row) is NOC. Let $C'$ be any odd cycle of $G'$. If $z \notin V(C')$, then $C'$ is an odd cycle in $G[U] \subseteq G$, so $\epsilon_{G'}(C') = \epsilon_G(C') = -1$. If $z \in V(C')$, write $C' = z w_1 P w_2 z$, where $w_1, w_2 \in U$ are distinct and $P$ is a simple $(w_1, w_2)$-path in $G[U]$ of odd length $2m - 1$ (so $|C'| = 2m + 1$). For each $i \in \{1, 2\}$, since $(A'_0)_{z, w_i} = a_{x, w_i} - a_{y, w_i} \ne 0$, the sign $\sigma'(z w_i)$ must equal either $\sigma(x w_i)$ (with $a_{x, w_i} \ne 0$) or $-\sigma(y w_i)$ (with $a_{y, w_i} \ne 0$).
\begin{itemize}
  \item If $\sigma'(z w_1) = \sigma(x w_1)$ and $\sigma'(z w_2) = \sigma(x w_2)$, then $C = x w_1 P w_2 x$ is a simple odd cycle of $G$ of the same length $2m + 1$ and the same edge signs as $C'$, so $\epsilon_{G'}(C') = \epsilon_G(C) = -1$.
  \item If $\sigma'(z w_1) = -\sigma(y w_1)$ and $\sigma'(z w_2) = -\sigma(y w_2)$, then $C = y w_1 P w_2 y$ is a simple odd cycle of $G$ of length $2m + 1$ whose sign product equals $(-1)^2$ times that of $C'$, so $\epsilon_{G'}(C') = \epsilon_G(C) = -1$.
  \item Otherwise, without loss of generality $\sigma'(z w_1) = \sigma(x w_1)$ and $\sigma'(z w_2) = -\sigma(y w_2)$. Because $V(P) \subseteq U$ is disjoint from $\{v, x, y\}$, the walk $C = x w_1 P w_2 y v x$ is a simple odd cycle of $G$ of length $(2m + 1) + 2 = 2(m + 1) + 1$. Using $\sigma(yv) = \sigma(vx) = +1$, we compute
  \begin{align*}
    \epsilon_G(C)
    &= (-1)^{m+1} \sigma(x w_1) \Bigl(\prod_{e \in P} \sigma(e)\Bigr) \sigma(y w_2)(+1)(+1) \\
    &= (-1)^m \sigma'(z w_1) \Bigl(\prod_{e \in P} \sigma'(e)\Bigr) \sigma'(z w_2) = \epsilon_{G'}(C').
  \end{align*}
  Hence $\epsilon_{G'}(C') = \epsilon_G(C) = -1$.
\end{itemize}
Thus $A'_0$ is NOC of order $n - 2$, and the induction hypothesis gives $s(A) \le s(A'_0) \le 0$.
\end{proof}

\subsection{SPQR-tree decomposition}

By Lemmas~\ref{lem:cutvertex} and~\ref{lem:deg2}, any minimum-order counterexample to Theorem~\ref{thm:main} is $2$-connected with $\delta(G) \ge 3$. If $G$ is not $3$-connected, we decompose it into its triconnected components via Tutte's decomposition, represented as an SPQR tree $\mathcal{T}$~\cite{Tut66,DT96} (with $Q$-nodes omitted). Each node $\mu$ of $\mathcal{T}$ carries a multigraph skeleton of type $S$ (a cycle of length at least $3$), $P$ (at least three parallel edges between two poles), or $R$ (a $3$-connected simple graph), and tree edges of $\mathcal{T}$ correspond to pairs of virtual edges glued along $2$-vertex cuts $\{x, y\}$ of $G$. As we verify in Section~\ref{sec:proof}, simplicity and $\delta(G) \ge 3$ force every leaf of $\mathcal{T}$ to be an $R$-node whose skeleton is a $3$-connected simple graph $T = H + xy$ with a single virtual edge $xy$.

This reduces the structural side of Theorem~\ref{thm:main} to understanding $3$-connected NOC signed graphs. We turn to that classification next.

% ============================================================================
\section{Classification of \texorpdfstring{$3$}{3}-Connected Non-Bipartite NOC Graphs}\label{sec:3conn}
% ============================================================================

For unsigned graphs, Bai, Grzesik, Li, and Prorok~\cite[Lemma~5.4]{BGLP26} (and independently Lin, Wang, and Zhou~\cite{LWZ25}) showed that every $3$-connected non-bipartite graph contains cycles of both odd lengths modulo $4$ unless it is $K_4$ or the Petersen graph. In this section we prove the signed-graph generalization needed for Theorem~\ref{thm:main}.

\begin{theorem}[classification of $3$-connected non-bipartite NOC graphs]\label{thm:3conn}
Let $(G, \sigma)$ be a $3$-connected non-bipartite signed graph. If $(G, \sigma)$ is NOC, then $(G, \sigma)$ is switching-equivalent to either $(K_4, +1)$ or $(\mathrm{Petersen}, -1)$.
\end{theorem}

\noindent\textbf{The parity mechanism.}
Throughout this section we write $w(e) = w_\sigma(e) \in \{1, 3\} \subset \mathbb{Z}/4\mathbb{Z}$. For a cycle $C$ with a fixed orientation and vertices $u, v \in V(C)$, let $\vecC{C}[u, v]$ and $\revC{C}[u, v]$ denote the two $(u, v)$-arcs of $C$. Because $w(e)$ is always odd, $2w(P) \equiv 2|P| \pmod 4$ for every path $P$ regardless of edge signs. Consequently, for any even cycle $C_1$ and any $u, v \in V(C_1)$,
\begin{equation}\label{eq:parity}
  w(\vecC{C}_1[u, v]) - w(\revC{C}_1[u, v]) \equiv w(C_1) - 2w(\revC{C}_1[u, v]) \equiv w(C_1) + 2d_{C_1}(u, v) \pmod 4.
\end{equation}
Thus the two arcs of $C_1$ between $u$ and $v$ have mod-$4$ weights differing by $2 \pmod 4$ whenever either
\begin{itemize}
  \item $w(C_1) \equiv 0 \pmod 4$ and $d_{C_1}(u, v)$ is odd, or
  \item $w(C_1) \equiv 2 \pmod 4$ and $d_{C_1}(u, v)$ is even.
\end{itemize}
Whenever two such vertices $u, v \in V(C_1)$ can be joined outside $C_1$ by a path $R$ whose length has the opposite parity to $d_{C_1}(u, v)$, the two odd cycles $R \cup \vecC{C}_1[u, v]$ and $R \cup \revC{C}_1[u, v]$ have weights differing by $2 \pmod 4$, so one of them must have weight $1 \pmod 4$ and violate NOC.

A similar observation forbids chords of even distance on any even cycle $C_1$ with $w(C_1) \equiv 2 \pmod 4$: if $u_0 v_0$ is a chord with $d_{C_1}(u_0, v_0)$ even, then the two odd cycles $D_1 = \vecC{C}_1[u_0, v_0] + u_0 v_0$ and $D_2 = \revC{C}_1[u_0, v_0] + u_0 v_0$ satisfy
\[
  w(D_1) + w(D_2) = w(C_1) + 2w(u_0 v_0) \equiv 2 + 2 \equiv 0 \pmod 4,
\]
so one of $D_1, D_2$ has weight $1 \pmod 4$, violating NOC.

To exploit this mechanism inside a $3$-connected graph, we use the maximal even-cycle lemma of Gao, Li, Ma, and Xie~\cite{GLMX24} (see also~\cite[Lemma~5.1]{BGLP26}).

\begin{lemma}[Gao--Li--Ma--Xie~\cite{GLMX24}]\label{lem:GLMX}
Let $G$ be a $3$-connected graph and $D$ a connected subgraph of $G$ such that $G - V(D)$ contains an even cycle. If $C$ is an even cycle in $G - V(D)$ chosen so that the component of $G - V(C)$ containing $D$ is maximal, then $G - V(C)$ is connected. Moreover, either $G[V(C)] \cong K_4$, or $C$ has at most one chord; and if $u_0 v_0$ is a chord of $C$, then the distance $d_C(u_0, v_0)$ along $C$ is even.
\end{lemma}

We first record the basic obstruction created when an even cycle and an odd cycle are linked by paths.

\begin{lemma}[even-odd cycle linkage]\label{lem:evodd}
Let $(G, \sigma)$ be a signed graph, $C_1$ an even cycle, and $C_2$ an odd cycle. Suppose that either:
\begin{enumerate}
  \item[(a)] $C_1$ and $C_2$ are vertex-disjoint and there are two vertex-disjoint paths $P^1, P^2$ from $C_1$ to $C_2$, internally disjoint from $C_1 \cup C_2$, with origins $u_1, u_2 \in V(C_1)$, or
  \item[(b)] $V(C_1) \cap V(C_2) = \{u_1\}$ and there is a path $P$ from $u_2 \in V(C_1) \setminus \{u_1\}$ to $V(C_2) \setminus \{u_1\}$ internally disjoint from $C_1 \cup C_2$.
\end{enumerate}
If $w(C_1) + 2d_{C_1}(u_1, u_2) \equiv 2 \pmod 4$, then $(G, \sigma)$ is not NOC.
\end{lemma}

\begin{proof}
In case~(a), let $v_1, v_2 \in V(C_2)$ be the endpoints of $P^1, P^2$ on $C_2$. Because $|C_2|$ is odd, the two $(u_1, u_2)$-paths $R_1 = P^1 \cup \vecC{C}_2[v_1, v_2] \cup P^2$ and $R_2 = P^1 \cup \revC{C}_2[v_1, v_2] \cup P^2$ are internally disjoint from $C_1$ and have opposite parities. Choose $R \in \{R_1, R_2\}$ so that $|R| + d_{C_1}(u_1, u_2)$ is odd. Then $R \cup \vecC{C}_1[u_1, u_2]$ and $R \cup \revC{C}_1[u_1, u_2]$ are two simple odd cycles whose mod-$4$ weights differ by $w(\vecC{C}_1[u_1, u_2]) - w(\revC{C}_1[u_1, u_2]) \equiv 2 \pmod 4$ by~\eqref{eq:parity}. One of them therefore has weight $1 \pmod 4$, so $(G, \sigma)$ is not NOC. Case~(b) is identical, taking $R_1, R_2$ to be $P$ followed by the two arcs of $C_2$ from the endpoint of $P$ on $C_2$ to $u_1$.
\end{proof}

Combining Lemmas~\ref{lem:GLMX} and~\ref{lem:evodd} imposes severe restrictions on pairs of cycles in a $3$-connected NOC signed graph: an even cycle and an odd cycle can neither be disjoint nor share a single vertex, and two disjoint odd cycles force the graph to be $(\mathrm{Petersen}, -1)$.

\begin{lemma}[cycle-pair obstructions]\label{lem:cyclepairs}
Let $(G, \sigma)$ be a $3$-connected NOC signed graph, and let $C_1, C_2$ be two cycles of $G$.
\begin{enumerate}
  \item[(1)] $C_1$ and $C_2$ cannot be vertex-disjoint with $C_1$ even and $C_2$ odd.
  \item[(2)] $C_1$ and $C_2$ cannot intersect in exactly one vertex with $C_1$ even and $C_2$ odd.
  \item[(3)] If $C_1$ and $C_2$ are vertex-disjoint and both odd, then $(G, \sigma)$ is switching-equivalent to $(\mathrm{Petersen}, -1)$.
\end{enumerate}
\end{lemma}

\begin{proof}
\emph{Proof of (1).} Suppose $C_1$ is even, $C_2$ is odd, and $V(C_1) \cap V(C_2) = \emptyset$. Applying Lemma~\ref{lem:GLMX} with $D = C_2$, choose an even cycle $C_1$ in $G - V(C_2)$ so that $H = G - V(C_1)$ is connected and either $G[V(C_1)] \cong K_4$ or $C_1$ has at most one chord, which joins two vertices at even distance on $C_1$.

If $G[V(C_1)] \cong K_4$, every triangle of $G[V(C_1)]$ has weight $3 \pmod 4$, so switching makes all six edges of $G[V(C_1)]$ positive. By $3$-connectivity, there exist two vertex-disjoint paths from $V(C_1)$ to $C_2$, internally disjoint from $C_1 \cup C_2$, with distinct origins $u_1, u_2 \in V(C_1)$. Inside $G[V(C_1)] \cong (K_4, +1)$ there are both a $1$-path and a $3$-path between $u_1$ and $u_2$ (weights $1$ and $3 \pmod 4$). Combining them with an even $(u_1, u_2)$-path routed through $C_2$ yields two odd cycles of weights differing by $2 \pmod 4$, contradicting NOC.

Thus $C_1$ has at most one chord (of even distance on $C_1$). Let $B$ be the block of $H$ containing $C_2$ (so $|V(B)| \ge |V(C_2)| \ge 3$). For each $v \in V(B)$, let $H_v$ be the component of $H - E(B)$ containing $v$ (the \emph{$B$-branch} at $v$). As in~\cite[Lemma~5.3]{BGLP26}, two vertices $y_1, y_2 \in V(H)$ admit vertex-disjoint paths in $H$ to $C_2$ if and only if they lie in distinct $B$-branches.

Suppose first that $w(C_1) \equiv 0 \pmod 4$. By Lemma~\ref{lem:evodd}(a), any two vertices of $C_1$ at odd distance that both have neighbors in $H$ must have all their neighbors in $H$ contained in the same $B$-branch. Let $(X, Y)$ be the bipartition of $C_1$. If both $X$ and $Y$ contain vertices with neighbors in $H$, then all of $N_H(V(C_1))$ lies in a single $B$-branch $H_v$, making $v$ a cut-vertex of $G$ separating $V(C_1)$ from $V(B) \setminus \{v\}$, impossible. Hence one part, say $Y$, satisfies $N_H(Y) = \emptyset$. Since $\delta(G) \ge 3$ and $C_1$ has at most one chord, at most two vertices of $Y$ can have degree at least $3$ in $G$, forcing $|Y| = 2$, $|C_1| = 4$, and $Y$ joined by the unique chord of $C_1$. Then $X$ is a $2$-vertex cut separating $Y$ from $H$, contradicting $3$-connectivity.

Suppose now that $w(C_1) \equiv 2 \pmod 4$. As observed above, $C_1$ cannot have a chord of even distance, so $C_1$ is chordless and every vertex of $C_1$ has a neighbor in $H$. By Lemma~\ref{lem:evodd}(a), any two distinct vertices of $C_1$ in the same part ($X$ or $Y$) have all their neighbors in $H$ contained in a single $B$-branch. Since $|X| \ge 2$ and $|Y| \ge 2$, we have $N_H(X) \subseteq V(H_{v_1})$ and $N_H(Y) \subseteq V(H_{v_2})$ for some $v_1, v_2 \in V(B)$. Because $|V(B)| \ge 3$, $\{v_1, v_2\}$ is a vertex cut of $G$ of size at most $2$ separating $C_1$ from $V(B) \setminus \{v_1, v_2\}$, contradicting $3$-connectivity.

\medskip
\emph{Proof of (2).} Suppose $V(C_1) \cap V(C_2) = \{u\}$ with $C_1$ even and $C_2$ odd. Applying Lemma~\ref{lem:GLMX} with $D = C_2 - u$, choose an even cycle $C_1$ in $G - V(C_2 - u)$ so that $H = G - V(C_1)$ is connected and either $G[V(C_1)] \cong K_4$ or $C_1$ has at most one chord (of even distance on $C_1$). If $u \notin V(C_1)$, then $C_1$ and $C_2$ are vertex-disjoint, contradicting~(1); hence $V(C_1) \cap V(C_2) = \{u\}$. The case $G[V(C_1)] \cong K_4$ is excluded just as in~(1) using a single path from $C_1 - u$ to $C_2 - u$ avoiding $u$.

If $w(C_1) \equiv 0 \pmod 4$, then by Lemma~\ref{lem:evodd}(b) no vertex of $C_1$ at odd distance from $u$ has a neighbor in $H$. Since $C_1$ has at most one chord and $\delta(G) \ge 3$, there can be at most two vertices at odd distance from $u$, forcing $|C_1| = 4$; writing $C_1 = u a v b u$, neither $a$ nor $b$ has a neighbor in $H$, so $\{u, v\}$ is a $2$-vertex cut of $G$, a contradiction.

If $w(C_1) \equiv 2 \pmod 4$, then $C_1$ is chordless, so every vertex of $C_1 - u$ has a neighbor in $H$. In particular, a vertex $v \in V(C_1)$ at distance $2$ from $u$ has a neighbor in $H$, yielding a path from $v$ to $C_2 - u$ avoiding $u$ and contradicting Lemma~\ref{lem:evodd}(b).

\medskip
\emph{Proof of (3).} Choose vertex-disjoint odd cycles $C_1, C_2$ minimizing $|C_1| + |C_2|$; then both $C_1$ and $C_2$ are chordless. If $G - (V(C_1) \cup V(C_2))$ had a component $Q$, then $3$-connectivity would give $|N_{C_1 \cup C_2}(Q)| \ge 3$, so $Q$ would have at least two neighbors on some $C_i$. A path through $Q$ between two neighbors on $C_i$ forms a $\Theta$-graph with $C_i$, which contains an even cycle disjoint from $C_{3-i}$, contradicting~(1). Hence $V(G) = V(C_1) \cup V(C_2)$. Similarly, if a vertex $x \in V(C_1)$ had two neighbors in $C_2$, then $G[V(C_2) \cup \{x\}]$ would contain a $\Theta$-graph and hence an even cycle meeting $C_1$ only at $x$, contradicting~(2). Therefore every vertex of $C_1$ has exactly one neighbor in $C_2$ and vice versa, so $|C_1| = |C_2| = L$ (odd, $L \ge 3$) and the edges between $C_1$ and $C_2$ form a perfect matching $M$. For $u \in V(C_1)$, denote its matching neighbor in $C_2$ by $u'$.

We claim that $d_{C_2}(u'_1, u'_2) \ge 2$ for every edge $u_1 u_2 \in E(C_1)$. Suppose to the contrary that $u'_1 u'_2 \in E(C_2)$ for some $u_1 u_2 \in E(C_1)$, forming the $4$-cycle $Q = u_1 u_2 u'_2 u'_1 u_1$. The odd cycle $C'_1 = (C_1 - u_1 u_2) \cup u_1 u'_1 u'_2 u_2$ satisfies
\[
  w(C'_1) \equiv w(C_1) + w(Q) - 2w(u_1 u_2) \equiv 3 + w(Q) + 2 \equiv w(Q) + 1 \pmod 4,
\]
so NOC forces $w(Q) \equiv 2 \pmod 4$.
\begin{itemize}
  \item If $L = 3$, then every edge of $C_1$ maps to an edge of $C_2$, yielding three $4$-cycles $Q_1, Q_2, Q_3$ around the triangular prism, each with $w(Q_i) \equiv 2 \pmod 4$. Summing their weights counts $C_1$, $C_2$, and each of the three matching edges twice:
  \[
    \sum_{i=1}^3 w(Q_i) \equiv w(C_1) + w(C_2) + 2 \cdot 3 \equiv 3 + 3 + 2 \equiv 0 \pmod 4,
  \]
  contradicting $2 + 2 + 2 \equiv 2 \pmod 4$.
  \item If $L \ge 5$, consider $H' = G - V(Q)$. If $H'$ contained an odd cycle, it would be disjoint from the even cycle $Q$, contradicting~(1). Thus $H'$ is bipartite. Let $P_1 = C_1 - u_1 u_2$ and $P_2 = C_2 - u'_1 u'_2$; both are paths of even length $L - 1 \ge 4$. Pick any vertex $x \in V(P_1) \setminus \{u_1, u_2\}$ and its matching neighbor $x' \in V(P_2) \setminus \{u'_1, u'_2\}$, and form the two cycles
  \[
    D_1 = P_1[u_1, x] x x' P_2[x', u'_1] u'_1 u_1, \qquad D_2 = P_1[u_2, x] x x' P_2[x', u'_2] u'_2 u_2.
  \]
  Because $L$ is odd, each of the paths $P_1 - \{u_1, u_2\}$ and $P_2 - \{u'_1, u'_2\}$ has $L - 2$ (odd) vertices, so its unique bipartition consists of a larger class of $(L - 1)/2$ vertices (at odd distance from $u_1$ on $P_1$, respectively from $u'_1$ on $P_2$) and a smaller class of $(L - 3)/2$ vertices (at even distance). Since $H' = (P_1 - \{u_1, u_2\}) \cup (P_2 - \{u'_1, u'_2\}) \cup (M \setminus \{u_1 u'_1, u_2 u'_2\})$ is bipartite, the perfect matching $M \setminus \{u_1 u'_1, u_2 u'_2\}$ must pair the two larger classes with each other, forcing $|P_1[u_1, x]| \equiv |P_2[u'_1, x']| \pmod 2$. Hence $|D_1| = |P_1[u_1, x]| + |P_2[x', u'_1]| + 2$ is even, and since $|P_1| = |P_2| = L - 1$ is even, $|D_2|$ is also even.

  Moreover, $D_1 \cup D_2$ traverses $P_1, P_2, u_1 u'_1, u_2 u'_2$ once and $xx'$ twice, so
  \begin{align*}
    w(D_1) + w(D_2) &\equiv w(P_1) + w(P_2) + w(u_1 u'_1) + w(u_2 u'_2) + 2w(xx') \\
    &\equiv (3 - w(u_1 u_2)) + (3 - w(u'_1 u'_2)) + (2 - w(u_1 u_2) - w(u'_1 u'_2)) + 2 \\
    &\equiv 10 - 2w(u_1 u_2) - 2w(u'_1 u'_2) \equiv 10 - 2 - 2 \equiv 2 \pmod 4.
  \end{align*}
  Thus $\{w(D_1), w(D_2)\} = \{0, 2\} \pmod 4$. Now $C_1 \triangle D_1$ and $C_1 \triangle D_2$ are simple odd cycles in $G$: $C_1 \triangle D_1 = P_1[x, u_2] u_2 u_1 u'_1 P_2[u'_1, x'] x' x$ shares $P_1[u_1, x]$ with $C_1$, and $C_1 \triangle D_2$ shares $P_1[x, u_2]$ with $C_1$. Since $|P_1| = L - 1$ is even, $d_1 = |P_1[u_1, x]|$ and $d_2 = |P_1[x, u_2]|$ have the same parity. If $d_1, d_2$ are even, pick $j \in \{1, 2\}$ with $w(D_j) \equiv 2 \pmod 4$; then $w(C_1 \triangle D_j) \equiv 3 + 2 - 0 \equiv 1 \pmod 4$. If $d_1, d_2$ are odd, pick $j \in \{1, 2\}$ with $w(D_j) \equiv 0 \pmod 4$; then $w(C_1 \triangle D_j) \equiv 3 + 0 - 2 \equiv 1 \pmod 4$. Either case contradicts NOC.
\end{itemize}
This proves $d_{C_2}(u'_1, u'_2) \ge 2$ for every edge $u_1 u_2 \in E(C_1)$.

In particular, $L \ne 3$. If $L = 5$, every adjacent pair in $C_1$ maps to vertices at distance $2$ in $C_2$, so $G$ is the Petersen graph. Because the $5$-cycles of the Petersen graph generate its cycle space over $\mathbb{F}_2$ (every $6$-cycle is the symmetric difference of two $5$-cycles sharing a $2$-path, and every chordless cycle of length $8$ or $9$ is a sum of $5$- and $6$-cycles), the NOC condition $\prod_{e \in C_5} \sigma(e) = -1$ on every $5$-cycle uniquely determines the sign product on every cycle of $G$. Hence $(G, \sigma)$ is switching-equivalent to $(\mathrm{Petersen}, -1)$.

Finally, suppose $L \ge 7$. Choose an edge $u_1 u_2 \in E(C_1)$ minimizing $d_{C_2}(u'_1, u'_2) \ge 2$, and orient $C_2$ so that $\vecC{C}_2[u'_1, u'_2]$ is even and $\revC{C}_2[u'_1, u'_2]$ is odd. Since $d_{C_2}(u'_1, u'_2) \ge 2$, the odd arc $\revC{C}_2[u'_1, u'_2]$ has length at least $3$, so it contains an edge $v'_1 v'_2$ with $\{v'_1, v'_2\} \cap \{u'_1, u'_2\} = \emptyset$. Let $v_1, v_2 \in V(C_1) \setminus \{u_1, u_2\}$ be their matching neighbors, and let $P$ be the $(v_1, v_2)$-path in $C_1 - \{u_1, u_2\}$. Then $C'_1 = u_1 u'_1 \vecC{C}_2[u'_1, u'_2] u'_2 u_2 u_1$ is an odd cycle, and $C'_2 = v'_1 v_1 P v_2 v'_2 v'_1$ is a cycle disjoint from $C'_1$. By~(1), $C'_2$ is odd. By minimality of $|C_1| + |C_2|$, we must have $V(G) = V(C'_1) \cup V(C'_2)$. This forces $|\revC{C}_2[u'_1, u'_2]| = 3$ (so $d_{C_2}(u'_1, u'_2) = 3$, as $L \ge 7$) and forces $v_1, v_2$ to be the two neighbors of $\{u_1, u_2\}$ on $C_1$. Hence either $u_1 v_1 \in E(C_1)$ or $u_1 v_2 \in E(C_1)$, giving an edge of $C_1$ whose matching endpoints have distance $1$ or $2$ on $C_2$, contradicting either $d_{C_2} \ge 2$ or the minimality of $d_{C_2}(u'_1, u'_2) = 3$.
\end{proof}

We now have all the ingredients needed to prove Theorem~\ref{thm:3conn}.

\begin{proof}[Proof of Theorem~\ref{thm:3conn}]
Let $C$ be a shortest odd cycle of $G$; then $C$ is chordless. If $G - V(C)$ contains a cycle $C'$, then Lemma~\ref{lem:cyclepairs}(1) rules out $C'$ even and Lemma~\ref{lem:cyclepairs}(3) forces $(G, \sigma)$ to be switching-equivalent to $(\mathrm{Petersen}, -1)$. Assume therefore that $G - V(C)$ is a forest, and let $T$ be a largest tree component of $G - V(C)$.

\emph{Case 1: $|C| = 3$, say $V(C) = \{v_1, v_2, v_3\}$.} If $T$ is a single vertex $u$, then $N_G(u) = V(C)$ by $3$-connectivity. If $G - V(C) = T$, then $G \cong K_4$, and since all triangles have $\epsilon = -1$ (sign product $+1$) and generate the cycle space of $K_4$, $(G, \sigma)$ is switching-equivalent to $(K_4, +1)$. If $G - V(C)$ has another component $T'$, then by maximality $V(T') = \{u'\}$ with $N_G(u') = V(C)$. Every triangle in $G[V(C) \cup \{u, u'\}]$ has sign product $+1$, so switching makes all edges of $G[V(C) \cup \{u, u'\}]$ positive; this produces the $5$-cycle $u v_1 u' v_2 v_3 u$ with $\epsilon = +1$, contradicting NOC.

Now suppose $|V(T)| \ge 2$, and let $u_1, u_2$ be two leaves of $T$ joined by the path $P$ in $T$. By $3$-connectivity, each $u_i$ has at least two neighbors in $C$, so $u_1, u_2$ share a neighbor $v_1 \in V(C)$. If $|P|$ is even, then $v_1 u_1 P u_2 v_1$ is an even cycle meeting $C$ only at $v_1$, contradicting Lemma~\ref{lem:cyclepairs}(2). Thus $|P|$ is odd. Let $v_2 \ne v_1$ be a second neighbor of $u_2$ in $C$, and consider the two odd cycles $C_a = v_1 u_1 P u_2 v_1$ and $C_b = v_1 u_1 P u_2 v_2 v_3 v_1$. Their symmetric difference is the $4$-cycle $Q = u_2 v_1 v_3 v_2 u_2$, whose chord $v_1 v_2$ splits $Q$ into the two triangles $C = v_1 v_2 v_3 v_1$ and $C' = u_2 v_1 v_2 u_2$. Since both triangles have weight $3 \pmod 4$,
\[
  w(Q) = w(C) + w(C') - 2w(v_1 v_2) \equiv 3 + 3 - 2 \equiv 0 \pmod 4.
\]
Hence $w(C_b) - w(C_a) \equiv w(Q) - 2w(u_2 v_1) \equiv 0 - 2 \equiv 2 \pmod 4$, so one of $C_a, C_b$ has weight $1 \pmod 4$, contradicting NOC.

\emph{Case 2: $|C| \ge 5$.} Because $C$ is a shortest odd cycle, any two neighbors $v_1, v_2 \in V(C)$ of a vertex $u \in V(T)$ satisfy $d_C(v_1, v_2) = 2$ (otherwise $u v_1$, $u v_2$, and the shorter arc of $C$ or its complement form a shorter odd cycle). In particular, every vertex of $T$ has at most two neighbors in $C$, so $|V(T)| \ge 2$. Moreover, for any $u \in V(T)$ with two neighbors $v_1, v_2 \in V(C)$ and common neighbor $v_0 \in V(C)$, let $Q_u = u v_1 v_0 v_2 u$ be the resulting $4$-cycle. Replacing $v_1 v_0 v_2$ in $C$ by $v_1 u v_2$ yields an odd cycle $C'$ with $w(C') - w(C) \equiv w(Q_u) \pmod 4$, so NOC forces $w(Q_u) \equiv 0 \pmod 4$.

Let $u_1$ be a leaf of $T$; by $3$-connectivity, $u_1$ has two neighbors $v_1, v_2 \in V(C)$ with common neighbor $v_0 \in V(C)$. If some vertex $v_3 \in V(C) \setminus \{v_0, v_1, v_2\}$ has a neighbor $u_3 \in V(T)$, let $P$ be the $(u_1, u_3)$-path in $T$, and assume $v_1, v_0, v_2, v_3$ appear in this order along $C$. Exactly one of the two cycles $\vecC{C}[v_1, v_3] v_3 u_3 P u_1 v_1$ and $\revC{C}[v_1, v_3] v_3 u_3 P u_1 v_1$ is odd; call it $C_1$. Replacing the path $v_1 v_0 v_2$ and edge $u_1 v_1$ by $u_1 v_2$ (if $C_1$ uses $\vecC{C}[v_1, v_3]$), or adding $v_1 v_0 v_2$ and swapping $u_1 v_1$ for $u_1 v_2$ (if $C_1$ uses $\revC{C}[v_1, v_3]$), produces a second odd cycle $C'_1$ with $C_1 \triangle C'_1 = Q_{u_1}$ and $|E(C_1) \cap E(Q_{u_1})| \in \{1, 3\}$. Since $w(Q_{u_1}) \equiv 0 \pmod 4$ and $2w(e) \equiv 2 \pmod 4$ for every edge $e$, we obtain
\[
  w(C_1) - w(C'_1) \equiv w(Q_{u_1}) - 2 \equiv 2 \pmod 4,
\]
contradicting NOC.

Therefore $N_C(T) \subseteq \{v_0, v_1, v_2\}$. Pick a second leaf $u_2 \ne u_1$ of $T$, and let $P$ be the $(u_1, u_2)$-path in $T$. Since $u_2$ has two neighbors in $C$ at distance $2$, we must have $N_C(u_2) = \{v_1, v_2\}$. If $|P|$ is even, then $v_1 u_1 P u_2 v_1$ is an even cycle meeting $C$ only at $v_1$, contradicting Lemma~\ref{lem:cyclepairs}(2). If $|P|$ is odd, then $C_1 = v_1 u_1 P u_2 v_1$ and $C_2 = v_1 u_1 P u_2 v_2 v_0 v_1$ are two odd cycles with $C_1 \triangle C_2 = Q_{u_2}$. Since $w(Q_{u_2}) \equiv 0 \pmod 4$, we get $w(C_2) - w(C_1) \equiv w(Q_{u_2}) - 2w(u_2 v_1) \equiv 2 \pmod 4$, contradicting NOC.
\end{proof}

Theorem~\ref{thm:3conn} settles the $3$-connected case once we verify that all weighted $(K_4, +1)$ and $(\mathrm{Petersen}, -1)$ matrices have non-positive signature. When $G$ is only $2$-connected, however, the SPQR tree decomposes $G$ into $R$-node leaves $T = H + xy$ whose subgraphs $H$ might be bipartite. We eliminate that possibility next.

% ============================================================================
\section{Bipartite \texorpdfstring{$R$}{R}-Node Leaves and Signed Path Weights}\label{sec:bipartite}
% ============================================================================

Suppose $T = H + xy$ is a $3$-connected simple graph with $H = T - xy$ bipartite, appearing as a leaf $R$-node of the SPQR tree of a non-bipartite NOC graph $G$. If $H$ contains two $(x, y)$-paths whose mod-$4$ weights differ by $2 \pmod 4$, then the complement $G_2$ of $H$ cannot contain any $(x, y)$-path of the opposite parity (otherwise closing both paths in $H$ with that path in $G_2$ would yield two odd cycles differing by $2 \pmod 4$). We now prove that such a pair of $(x, y)$-paths in $H$ \emph{always} exists, for every signing $\sigma$ of $H$.

\begin{proposition}[mod-$4$ path-weight diversity in bipartite $R$-leaves]\label{prop:bipartite}
Let $T = H + xy$ be a $3$-connected simple graph such that $H = T - xy$ is bipartite, and let $\sigma$ be any signing of $H$. Then $H$ contains two $(x, y)$-paths $P_1, P_2$ such that $w_\sigma(P_1) - w_\sigma(P_2) \equiv 2 \pmod 4$.
\end{proposition}

\begin{proof}
Suppose for contradiction that every $(x, y)$-path in $H$ has the same mod-$4$ weight $r \in \mathbb{Z}/4\mathbb{Z}$.

\emph{Case 1: $x$ and $y$ lie in the same bipartition class of $H$.} Then $r \in \{0, 2\}$ and $T = H + xy$ is $3$-connected and non-bipartite. Extend $\sigma$ to $T$ by choosing $\sigma(xy) \in \{+1, -1\}$ so that $w_\sigma(xy) + r \equiv 3 \pmod 4$. Since $H$ is bipartite, every odd cycle of $T$ traverses the edge $xy$ and therefore consists of $xy$ together with an $(x, y)$-path $P$ in $H$. Every such cycle has weight $w_\sigma(xy) + w_\sigma(P) \equiv 3 \pmod 4$, so $(T, \sigma)$ is a $3$-connected non-bipartite NOC signed graph. By Theorem~\ref{thm:3conn}, $T \cong K_4$ or $T \cong \mathrm{Petersen}$. However, neither $K_4 - xy$ nor $\mathrm{Petersen} - xy$ is bipartite, a contradiction.

\emph{Case 2: $x$ and $y$ lie in opposite bipartition classes of $H$.} Then $T = H + xy$ is a $3$-connected bipartite graph. Every vertex $v \in V(T) \setminus \{x, y\}$ has degree at least $3$ in $T$ and can be adjacent to at most one of $x, y$ (because $xy \in E(T)$ and $T$ is bipartite), so $\delta(T - \{x, y\}) \ge 2$. Hence $T - \{x, y\}$ contains a cycle, which is necessarily even.

Apply Lemma~\ref{lem:GLMX} to $T$ with $D = \{x, y\}$ (the edge $xy$): let $C$ be an even cycle in $T - \{x, y\}$ maximizing the component $K$ of $T - V(C)$ containing $\{x, y\}$. Then $K = T - V(C)$ is connected. Since $T$ is bipartite, $T[V(C)] \not\cong K_4$ and $C$ has no chord of even distance; by Lemma~\ref{lem:GLMX}, $C$ has no chord of odd distance either. Thus $C$ is chordless, and every vertex of $C$ has a neighbor in $K$.

Let $B$ be the unique block of $K$ containing the edge $xy$, and for each $v \in V(B)$, let $K_v$ be the $B$-branch at $v$ (the component of $K - E(B)$ containing $v$). Whether $B$ is the single edge $xy$ or a $2$-connected subgraph containing $x$ and $y$, for any two distinct vertices $v_1, v_2 \in V(B)$ there exist two vertex-disjoint paths in $B - xy$ from $\{x, y\}$ to $\{v_1, v_2\}$ (any two vertex-disjoint paths in $B$ from $\{x, y\}$ to $\{v_1, v_2\}$ start at the distinct endpoints $x$ and $y$, and therefore automatically avoid the edge $xy$). Consequently, for any two vertices $u_1, u_2 \in V(C)$ whose neighbors in $K$ lie in distinct $B$-branches, there exist two vertex-disjoint paths $L_1, L_2$ in $H$ from $\{x, y\}$ to $\{u_1, u_2\}$ internally disjoint from $C$. Concatenating $L_1, L_2$ with $\vecC{C}[u_1, u_2]$ and $\revC{C}[u_1, u_2]$ yields two $(x, y)$-paths in $H$ whose mod-$4$ weights differ by $w_\sigma(C) + 2d_C(u_1, u_2) \pmod 4$ by~\eqref{eq:parity}. By assumption, this difference must vanish modulo $4$.
\begin{itemize}
  \item If $w_\sigma(C) \equiv 0 \pmod 4$, then no two vertices of $C$ at odd distance can have neighbors in distinct $B$-branches. Since $C$ is connected and every vertex of $C$ has a neighbor in $K$, all of $N_K(V(C))$ lies in a single $B$-branch $K_v$. Then $v$ is a cut-vertex of $T$ separating $C$ from $V(B) \setminus \{v\}$, contradicting $3$-connectivity.
  \item If $w_\sigma(C) \equiv 2 \pmod 4$, let $(X, Y)$ be the bipartition of $C$. No two distinct vertices in $X$ (which are at even distance on $C$) can attach to distinct $B$-branches, so $N_K(X) \subseteq V(K_{v_1})$ for some $v_1 \in V(B)$; similarly $N_K(Y) \subseteq V(K_{v_2})$ for some $v_2 \in V(B)$. Clearly $v_1 \ne v_2$, else $v_1$ is a cut-vertex of $T$. If $|V(B)| \ge 3$, then $\{v_1, v_2\}$ is a $2$-vertex cut of $T$ separating $C$ from $V(B) \setminus \{v_1, v_2\}$, contradicting $3$-connectivity. Hence $V(B) = \{v_1, v_2\} = \{x, y\}$, so $K_x$ and $K_y$ are vertex-disjoint with $N_K(X) \subseteq V(K_x)$ and $N_K(Y) \subseteq V(K_y)$ (after possibly swapping $X$ and $Y$).

  Since $|C| \ge 4$, choose four vertices $a_1, b_1, a_2, b_2$ appearing in cyclic order along $C$ with $a_1, a_2 \in X$ and $b_1, b_2 \in Y$. Let $A_1 = \vecC{C}[a_1, b_1]$, $A_2 = \vecC{C}[b_1, a_2]$, $A_3 = \vecC{C}[a_2, b_2]$, and $A_4 = \vecC{C}[b_2, a_1]$ be the four arcs partitioning $C$; each $A_i$ has odd length. Let $\revC{A}_2$ and $\revC{A}_4$ denote the reversals of $A_2$ (from $a_2$ to $b_1$) and $A_4$ (from $a_1$ to $b_2$). For $i = 1, 2$, choose a path $L_i$ from $x$ to $a_i$ whose vertices other than $a_i$ lie in $K_x$, and a path $R_i$ from $b_i$ to $y$ whose vertices other than $b_i$ lie in $K_y$. Because $K_x$, $K_y$, and $C$ are pairwise vertex-disjoint, the four walks
  \[
    P_1 = L_1 A_1 R_1, \qquad P_2 = L_2 \revC{A}_2 R_1, \qquad P_3 = L_2 A_3 R_2, \qquad P_4 = L_1 \revC{A}_4 R_2
  \]
  are simple $(x, y)$-paths in $H$. Their alternating weight sum cancels $L_1, L_2, R_1, R_2$:
  \begin{align*}
    \bigl(w_\sigma(P_1) + w_\sigma(P_3)\bigr) - \bigl(w_\sigma(P_2) + w_\sigma(P_4)\bigr)
    &\equiv w_\sigma(A_1) + w_\sigma(A_3) - w_\sigma(A_2) - w_\sigma(A_4) \\
    &\equiv w_\sigma(C) - 2\bigl(w_\sigma(A_2) + w_\sigma(A_4)\bigr) \pmod 4.
  \end{align*}
  Since $|A_2|$ and $|A_4|$ are both odd, $w_\sigma(A_2) + w_\sigma(A_4)$ is even, so $2(w_\sigma(A_2) + w_\sigma(A_4)) \equiv 0 \pmod 4$. Therefore
  \[
    \bigl(w_\sigma(P_1) + w_\sigma(P_3)\bigr) - \bigl(w_\sigma(P_2) + w_\sigma(P_4)\bigr) \equiv w_\sigma(C) \equiv 2 \pmod 4,
  \]
  contradicting $w_\sigma(P_1) \equiv w_\sigma(P_2) \equiv w_\sigma(P_3) \equiv w_\sigma(P_4) \equiv r \pmod 4$. \qedhere
\end{itemize}
\end{proof}

With Theorem~\ref{thm:3conn} and Proposition~\ref{prop:bipartite} in place, every $3$-connected non-bipartite NOC graph, and every leaf $R$-node of a non-bipartite $2$-connected NOC graph, arises from $(K_4, +1)$ or $(\mathrm{Petersen}, -1)$. All that remains is the linear algebra showing that both configurations can be eliminated without increasing the signature.

% ============================================================================
\section{Schur-Complement Reductions for \texorpdfstring{$K_4$}{K4} and Petersen Leaves}\label{sec:schur}
% ============================================================================

We now prove the two algebraic reduction lemmas handling $(K_4, +1)$ and $(\mathrm{Petersen}, -1)$, both as standalone $3$-connected graphs and as $R$-node leaves attached at a $2$-vertex cut $\{x, y\}$. In each leaf reduction, eliminating the internal vertices produces a $2 \times 2$ Schur-complement matrix $M_H$ on $\{x, y\}$ whose diagonal entries are strictly negative (hence removable via Weyl's monotonicity theorem) and whose off-diagonal entry has the \emph{exact} sign compatible with a path through $H$ of the same mod-$4$ weight—so replacing $H$ by a single weighted edge $xy$ preserves the NOC property.

\begin{lemma}[positive-weighted $K_4$ and diamond reduction]\label{lem:K4}
Let $A$ be a real symmetric zero-diagonal matrix with underlying NOC signed graph $(G, \sigma)$.
\begin{enumerate}
  \item[(1)] If $(G, \sigma)$ is switching-equivalent to $(K_4, +1)$, then $s(A) \le 0$.
  \item[(2)] Suppose $G = H \cup G_2$ with $V(H) \cap V(G_2) = \{x, y\}$, $xy \notin E(H)$, and $(H, \sigma|_H)$ switching-equivalent to $(K_4 - xy, +1)$. If Theorem~\ref{thm:main} holds for all NOC matrices of order $|V(G)| - 2$, then $s(A) \le 0$.
\end{enumerate}
\end{lemma}

\begin{proof}
For~(1), after switching we may assume $a_{ij} > 0$ for all $1 \le i < j \le 4$. Let $A_3$ be the leading $3 \times 3$ principal submatrix. Then $\tr(A_3) = 0$ and $\det(A_3) = 2 a_{12} a_{23} a_{13} > 0$, so the eigenvalues of $A_3$ satisfy $\lambda_1 > 0 > \lambda_2 \ge \lambda_3$; hence $n^-(A_3) = 2$. By Cauchy's interlacing theorem, $n^-(A) \ge n^-(A_3) = 2$, so $n^+(A) \le 4 - n^-(A) \le 2$ and $s(A) = n^+(A) - n^-(A) \le 0$.

For~(2), after switching on $V(G)$, we may assume all five edges of $H = K_4 - xy$ have positive weights. Write $V(H) \setminus \{x, y\} = \{u, v\}$, and set $a = a_{xu}$, $b = a_{xv}$, $c = a_{yu}$, $d = a_{yv}$, $e = a_{uv} > 0$. Ordering the vertices as $\{u, v\}, \{x, y\}, V(G_2) \setminus \{x, y\}$, the interior block is $A_{\mathrm{int}} = \begin{pmatrix} 0 & e \\ e & 0 \end{pmatrix}$, which is nonsingular with $s(A_{\mathrm{int}}) = 0$ and $A_{\mathrm{int}}^{-1} = \frac{1}{e}\begin{pmatrix} 0 & 1 \\ 1 & 0 \end{pmatrix}$. Its Schur complement on $\{x, y\}$ is
\[
  M_H = -\begin{pmatrix} a & b \\ c & d \end{pmatrix} \frac{1}{e}\begin{pmatrix} 0 & 1 \\ 1 & 0 \end{pmatrix} \begin{pmatrix} a & c \\ b & d \end{pmatrix}
  = \begin{pmatrix} -2ab/e & -(ad+bc)/e \\ -(ad+bc)/e & -2cd/e \end{pmatrix}.
\]
Thus $s(A) = s(A(G_2) + M_H \oplus 0)$. Since $-2ab/e < 0$ and $-2cd/e < 0$, we have
\[
  M_H \preceq \begin{pmatrix} 0 & -\mu \\ -\mu & 0 \end{pmatrix}, \qquad \mu = \frac{ad + bc}{e} > 0.
\]
Let $A'$ be the $(|V(G)| - 2) \times (|V(G)| - 2)$ zero-diagonal matrix obtained from $A(G_2)$ by subtracting $\mu$ from the $(x, y)$ and $(y, x)$ entries. By Weyl's monotonicity theorem, $s(A) \le s(A')$. To check that $(G(A'), \sigma')$ is NOC, note that only the entry $a'_{xy} = a_{xy} - \mu$ changes. If $a'_{xy} = 0$ or $\sgn(a'_{xy}) = \sgn(a_{xy})$ (with $a_{xy} \ne 0$), then $(G(A'), \sigma')$ is a signed subgraph of $G_2$, hence NOC. If $a'_{xy} < 0$, any odd cycle $C'$ of $G(A')$ using $xy$ lifts to a simple odd cycle $C$ of $G$ by replacing the negative edge $xy$ with the positive $3$-path $xuvy$ in $H$. Since $w_\sigma(xuvy) = 1 + 1 + 1 = 3 = w_{\sigma'}(xy)$, we have $\epsilon_G(C) = \epsilon_{G(A')}(C')$, so $(G(A'), \sigma')$ is NOC. By hypothesis, $s(A) \le s(A') \le 0$.
\end{proof}

\noindent\textbf{The $4 \times 4$ block decomposition of the Petersen interior.}
At first sight, eliminating the $8$-vertex interior $V(H) \setminus \{x, y\}$ of a weighted $(\mathrm{Petersen} - xy, -1)$ leaf looks forbidding: $A_{\mathrm{int}}$ is an $8 \times 8$ matrix with $14$ independent negative edge weights, and it can even be singular. The key observation is combinatorial. Label the Petersen graph with outer $5$-cycle $0\text{--}1\text{--}2\text{--}3\text{--}4\text{--}0$, inner $5$-cycle $5\text{--}7\text{--}9\text{--}6\text{--}8\text{--}5$, and spokes $i\text{--}(i+5)$, taking $x = 0$ and $y = 1$. The $8$ internal vertices split cleanly into two $4$-sets:
\begin{itemize}
  \item $U = \{2, 4, 5, 6\}$, which is an \emph{independent set} containing all neighbors of $0$ and $1$ in $H$ ($N_H(0) = \{4, 5\}$ and $N_H(1) = \{2, 6\}$);
  \item $W = \{3, 7, 8, 9\}$, which has \emph{no edges} to $\{0, 1\}$ and induces just two disjoint edges ($3\text{--}8$ and $7\text{--}9$).
\end{itemize}
With this partition, $A_{\mathrm{int}} = \begin{pmatrix} 0 & B \\ B^T & C \end{pmatrix}$ and the coupling matrix $R = \begin{pmatrix} R_U \\ 0 \end{pmatrix}$ is supported solely on $U$. Because $B$ is $4 \times 4$ with only two nonzero entries per row and column, $\det(B)$ is a difference of two monomials and every cofactor of $B$ has a \emph{fixed sign} for all positive weight magnitudes. This makes both the nonsingular case ($\det(B) \ne 0$) and the singular case ($\det(B) = 0$) completely transparent.

\begin{lemma}[negative-weighted Petersen and Petersen-leaf reduction]\label{lem:Petersen}
Let $A$ be a real symmetric zero-diagonal matrix with underlying NOC signed graph $(G, \sigma)$.
\begin{enumerate}
  \item[(1)] If $(G, \sigma)$ is switching-equivalent to $(\mathrm{Petersen}, -1)$, then $s(A) \le 0$.
  \item[(2)] Suppose $G = H \cup G_2$ with $V(H) \cap V(G_2) = \{x, y\}$, $xy \notin E(H)$, and $(H, \sigma|_H)$ switching-equivalent to $(\mathrm{Petersen} - xy, -1)$. If Theorem~\ref{thm:main} holds for all NOC matrices of order at most $|V(G)| - 8$, then $s(A) \le 0$.
\end{enumerate}
\end{lemma}

\begin{proof}
We treat parts~(1) and~(2) simultaneously, viewing $(\mathrm{Petersen}, -1)$ as the special case $G_2 = (\{x, y\}, \{xy\})$ with $V(G_2) \setminus \{x, y\} = \emptyset$. After switching on $V(G)$, we may assume all $14$ edges of $H = \mathrm{Petersen} - xy$ have strictly negative weights. With the vertex labeling above ($x = 0$, $y = 1$), denote the $14$ positive magnitudes $w_k > 0$ of the edge weights $-w_k$ of $H$ by
\begin{align*}
  w_0 &= |a_{12}|, & w_1 &= |a_{23}|, & w_2 &= |a_{34}|, & w_3 &= |a_{40}|, & w_4 &= |a_{57}|, & w_5 &= |a_{68}|, & w_6 &= |a_{79}|, \\
  w_7 &= |a_{85}|, & w_8 &= |a_{96}|, & w_9 &= |a_{05}|, & w_{10} &= |a_{16}|, & w_{11} &= |a_{27}|, & w_{12} &= |a_{38}|, & w_{13} &= |a_{49}|.
\end{align*}
Ordering $V(H) \setminus \{0, 1\} = U \cup W$ as $(2, 4, 5, 6;\; 3, 7, 8, 9)$, the interior matrix $A_{\mathrm{int}}$ and the $8 \times 2$ attachment matrix $R$ of $\{0, 1\}$ to $V(H) \setminus \{0, 1\}$ take the block forms
\[
  A_{\mathrm{int}} = \begin{pmatrix} 0 & B \\ B^T & C \end{pmatrix}, \qquad R = \begin{pmatrix} R_U \\ 0 \end{pmatrix},
\]
where
\[
  B = \begin{pmatrix}
    -w_1 & -w_{11} & 0 & 0 \\
    -w_2 & 0 & 0 & -w_{13} \\
    0 & -w_4 & -w_7 & 0 \\
    0 & 0 & -w_5 & -w_8
  \end{pmatrix}, \quad
  C = \begin{pmatrix}
    0 & 0 & -w_{12} & 0 \\
    0 & 0 & 0 & -w_6 \\
    -w_{12} & 0 & 0 & 0 \\
    0 & -w_6 & 0 & 0
  \end{pmatrix}, \quad
  R_U = \begin{pmatrix}
    0 & -w_0 \\
    -w_3 & 0 \\
    -w_9 & 0 \\
    0 & -w_{10}
  \end{pmatrix}.
\]
For any $v = (v_3, v_7, v_8, v_9)^T \in \mathbb{R}^4$, we have $v^T C v = -2w_{12} v_3 v_8 - 2w_6 v_7 v_9$. Expanding $\det(B)$ along the first row gives
\[
  \det(B) = w_1 w_4 w_5 w_{13} - w_2 w_7 w_8 w_{11}, \qquad \det(A_{\mathrm{int}}) = \det(B)^2.
\]
We analyze the nonsingular and singular cases separately.

\medskip
\emph{Case 1: $\det(B) \ne 0$.} Setting $P = \begin{pmatrix} I_4 & -\frac{1}{2}(B^T)^{-1}C \\ 0 & I_4 \end{pmatrix}$, we obtain $P^T A_{\mathrm{int}} P = \begin{pmatrix} 0 & B \\ B^T & 0 \end{pmatrix}$, whose spectrum is symmetric about $0$; hence $s(A_{\mathrm{int}}) = 0$. Moreover,
\[
  A_{\mathrm{int}}^{-1} = \begin{pmatrix} -(B^T)^{-1} C B^{-1} & (B^T)^{-1} \\ B^{-1} & 0 \end{pmatrix}.
\]
Because $R = \begin{pmatrix} R_U \\ 0 \end{pmatrix}$, the $2 \times 2$ Schur complement $M_H = -R^T A_{\mathrm{int}}^{-1} R$ on $\{0, 1\}$ simplifies to
\[
  M_H = (B^{-1} R_U)^T C (B^{-1} R_U) = \frac{1}{\det(B)^2} \begin{pmatrix} v^{(0)T} C v^{(0)} & v^{(0)T} C v^{(1)} \\ v^{(0)T} C v^{(1)} & v^{(1)T} C v^{(1)} \end{pmatrix},
\]
where $v^{(0)} = \adj(B) R_{U,0}$ and $v^{(1)} = \adj(B) R_{U,1}$ are the two columns of $\adj(B) R_U$. Direct evaluation of the $3 \times 3$ cofactors of $B$ yields
\begin{align*}
  v^{(0)}_3 &= -w_{11}(w_{13} w_5 w_9 + w_3 w_7 w_8), & v^{(0)}_7 &= w_1(w_{13} w_5 w_9 + w_3 w_7 w_8), \\
  v^{(0)}_8 &= -w_8(w_1 w_3 w_4 + w_{11} w_2 w_9), & v^{(0)}_9 &= w_5(w_1 w_3 w_4 + w_{11} w_2 w_9),
\end{align*}
and
\begin{align*}
  v^{(1)}_3 &= w_{13}(w_0 w_4 w_5 + w_{10} w_{11} w_7), & v^{(1)}_7 &= -w_7(w_0 w_2 w_8 + w_1 w_{10} w_{13}), \\
  v^{(1)}_8 &= w_4(w_0 w_2 w_8 + w_1 w_{10} w_{13}), & v^{(1)}_9 &= -w_2(w_0 w_4 w_5 + w_{10} w_{11} w_7).
\end{align*}
Because every $w_k > 0$, the sign patterns of $v^{(0)}$ and $v^{(1)}$ are strictly $(-, +, -, +)^T$ and $(+, -, +, -)^T$, respectively. Thus $v^{(0)}_3 v^{(0)}_8 > 0$, $v^{(0)}_7 v^{(0)}_9 > 0$, $v^{(1)}_3 v^{(1)}_8 > 0$, and $v^{(1)}_7 v^{(1)}_9 > 0$, while all four cross-products $v^{(0)}_3 v^{(1)}_8$, $v^{(0)}_8 v^{(1)}_3$, $v^{(0)}_7 v^{(1)}_9$, $v^{(0)}_9 v^{(1)}_7$ are strictly negative. Consequently:
\begin{align*}
  (M_H)_{00} &= \frac{-2w_{12} v^{(0)}_3 v^{(0)}_8 - 2w_6 v^{(0)}_7 v^{(0)}_9}{\det(B)^2} < 0, \\
  (M_H)_{11} &= \frac{-2w_{12} v^{(1)}_3 v^{(1)}_8 - 2w_6 v^{(1)}_7 v^{(1)}_9}{\det(B)^2} < 0, \\
  (M_H)_{01} &= \frac{-w_{12}\bigl(v^{(0)}_3 v^{(1)}_8 + v^{(0)}_8 v^{(1)}_3\bigr) - w_6\bigl(v^{(0)}_7 v^{(1)}_9 + v^{(0)}_9 v^{(1)}_7\bigr)}{\det(B)^2} > 0.
\end{align*}

In part~(1) (full negative-weighted Petersen), $s(A) = s(A_{\mathrm{int}}) + s\bigl(M_H + \begin{pmatrix} 0 & a_{01} \\ a_{01} & 0 \end{pmatrix}\bigr)$. The $2 \times 2$ matrix $M_H + \begin{pmatrix} 0 & a_{01} \\ a_{01} & 0 \end{pmatrix}$ has strictly negative trace $(M_H)_{00} + (M_H)_{11} < 0$, so it cannot have two positive eigenvalues; hence its signature is at most $0$, giving $s(A) \le 0$.

In part~(2), put $\mu = (M_H)_{01} > 0$. Since $(M_H)_{00} < 0$ and $(M_H)_{11} < 0$, we have $M_H \preceq \begin{pmatrix} 0 & \mu \\ \mu & 0 \end{pmatrix}$. Let $A'$ be the $(|V(G)| - 8) \times (|V(G)| - 8)$ zero-diagonal matrix obtained from $A(G_2)$ by adding $\mu$ to the $(x, y)$ and $(y, x)$ entries. By Weyl's monotonicity theorem, $s(A) = s(A(G_2) + M_H \oplus 0) \le s(A')$. If $a'_{xy} = a_{xy} + \mu$ is zero or has the same sign as $a_{xy}$ (with $a_{xy} \ne 0$), then $(G(A'), \sigma')$ is a signed subgraph of $G_2$, hence NOC. If $a'_{xy} > 0$, any odd cycle $C'$ of $G(A')$ using the positive edge $xy$ lifts to a simple odd cycle $C$ of $G$ by replacing $xy$ with the $7$-path $0\text{--}5\text{--}8\text{--}3\text{--}4\text{--}9\text{--}6\text{--}1$ in $H$. Since all $7$ edges of this path are negative, its mod-$4$ weight is $7 \cdot 3 \equiv 1 = w_{\sigma'}(xy) \pmod 4$, so $\epsilon_G(C) = \epsilon_{G(A')}(C')$. Thus $A'$ is NOC of order $|V(G)| - 8$, and by hypothesis $s(A) \le s(A') \le 0$.

\medskip
\emph{Case 2: $\det(B) = 0$.} Deleting row~$1$ and column~$1$ of $B$ leaves the $3 \times 3$ submatrix
\[
  \begin{pmatrix}
    0 & 0 & -w_{13} \\
    -w_4 & -w_7 & 0 \\
    0 & -w_5 & -w_8
  \end{pmatrix},
\]
whose determinant is $-w_4 w_5 w_{13} \ne 0$, so $\rank(B) = 3$. The unique nonzero kernel vectors $p \in \ker(B^T)$ and $q \in \ker(B)$ normalized by $p_2 = 1$ and $q_3 = 1$ are
\[
  p = \left(1,\, -\frac{w_1}{w_2},\, -\frac{w_{11}}{w_4},\, \frac{w_{11} w_7}{w_4 w_5}\right)^T, \qquad
  q = \left(1,\, -\frac{w_1}{w_{11}},\, \frac{w_1 w_4}{w_{11} w_7},\, -\frac{w_2}{w_{13}}\right)^T.
\]
Hence $q$ has sign pattern $(+, -, +, -)^T$, which gives
\[
  q^T C q = -2w_{12} q_3 q_8 - 2w_6 q_7 q_9 < 0.
\]
Extending $p$ to a basis of $\mathbb{R}^U$ and $q$ to a basis of $\mathbb{R}^W$ transforms $B$ into $\begin{pmatrix} B_1 & 0 \\ 0 & 0 \end{pmatrix}$ with $B_1$ nonsingular of size $3 \times 3$, and transforms $A_{\mathrm{int}}$ by congruence into
\[
  \begin{pmatrix}
    0 & 0 & B_1 & 0 \\
    0 & 0 & 0 & 0 \\
    B_1^T & 0 & C_{11} & c_{12} \\
    0 & 0 & c_{12}^T & q^T C q
  \end{pmatrix}.
\]
The vector $\tilde{p} = \begin{pmatrix} p \\ 0 \end{pmatrix} \in \mathbb{R}^8$ spans $\ker(A_{\mathrm{int}})$. Let $A_{\mathrm{int}}|_S$ be the remaining $7 \times 7$ principal block on the complementary $7$-dimensional subspace $S$ spanned by the other basis vectors (which is congruent to the orthogonal restriction $A_{\mathrm{int}}|_{\tilde{p}^\perp}$ since $\tilde{p} \in \ker(A_{\mathrm{int}})$). Its leading $6 \times 6$ submatrix $K = \begin{pmatrix} 0 & B_1 \\ B_1^T & C_{11} \end{pmatrix}$ is nonsingular with $s(K) = 0$ and its inverse has a $3 \times 3$ zero block in its bottom-right corner; thus the Schur complement of $K$ on $q$ is simply $q^T C q < 0$. Consequently, $A_{\mathrm{int}}|_S$ is nonsingular of size $7 \times 7$ with signature
\[
  s(A_{\mathrm{int}}|_S) = s(K) + s(q^T C q) = 0 + (-1) = -1.
\]

Now compute the inner products of $\tilde{p}$ with the two columns of $R$:
\begin{align*}
  \alpha &= p^T R_{U,0} = -w_3 p_4 - w_9 p_5 = \frac{w_1 w_3}{w_2} + \frac{w_{11} w_9}{w_4} > 0, \\
  \beta  &= p^T R_{U,1} = -w_0 p_2 - w_{10} p_6 = -w_0 - \frac{w_{10} w_{11} w_7}{w_4 w_5} < 0.
\end{align*}
Change basis on $\{x, y\}$ to $z_1 = x$ and $z_2 = (-\beta)x + \alpha y$ (invertible since $\alpha > 0$), so that $\tilde{p}^T R z_1 = \alpha \ne 0$ and $\tilde{p}^T R z_2 = -\beta\alpha + \alpha\beta = 0$. In the basis $\bigl(\tilde{p}, z_1;\; S;\; z_2, V(G_2) \setminus \{x, y\}\bigr)$, the pair $(\tilde{p}, z_1)$ forms a leading $2 \times 2$ block $\begin{pmatrix} 0 & \alpha \\ \alpha & 0 \end{pmatrix}$ (signature $0$) whose first row and column ($\tilde{p}$) have zero entries to all remaining basis vectors. Eliminating $(\tilde{p}, z_1)$ therefore adds zero to the remaining $(|V(G)| - 2) \times (|V(G)| - 2)$ block, which is
\[
  \begin{pmatrix}
    A_{\mathrm{int}}|_S & \tilde{r} & 0 \\
    \tilde{r}^T & -2\alpha\beta a_{xy} & \hat{r}^T \\
    0 & \hat{r} & A_2
  \end{pmatrix},
\]
where $A_2 = A(G_2 - \{x, y\})$ (empty in part~(1)) and $\hat{r} = (-\beta) r_{2,x} + \alpha r_{2,y}$. Eliminating the nonsingular $7 \times 7$ block $A_{\mathrm{int}}|_S$ (of signature $-1$) gives
\[
  s(A) = -1 + s(M), \qquad M = \begin{pmatrix} d & \hat{r}^T \\ \hat{r} & A_2 \end{pmatrix},
\]
where $d = -2\alpha\beta a_{xy} - \tilde{r}^T (A_{\mathrm{int}}|_S)^{-1} \tilde{r} \in \mathbb{R}$.

In part~(1) (and whenever $V(G_2) \setminus \{x, y\} = \emptyset$), $A_2$ is empty, so $M = (d)$ is $1 \times 1$, hence $s(M) \le 1$ and $s(A) = -1 + s(M) \le 0$. In part~(2) with $V(G_2) \setminus \{x, y\} \ne \emptyset$, $A_2 = A(G_2 - \{x, y\})$ is an induced NOC subgraph of order $|V(G)| - 10$, so $s(A_2) \le 0$ by hypothesis. If $\hat{r} \notin \im(A_2)$, pick $u \in \ker(A_2)$ with $u^T \hat{r} = 1$; extending $u$ to a basis of $\mathbb{R}^{V(G_2) \setminus \{x, y\}}$, the $2 \times 2$ block on $(z_2, u)$ is $\begin{pmatrix} d & 1 \\ 1 & 0 \end{pmatrix}$ (signature $0$), whose second row and column ($u$) have zero entries to the remaining basis vectors, so its Schur complement on the rest of $A_2$ is zero; hence $s(M) = s(A_2) \le 0$ and $s(A) \le -1$. If $\hat{r} \in \im(A_2)$, choose $y_2$ with $A_2 y_2 = \hat{r}$; block congruence gives $s(M) = s(A_2) + s(d - \hat{r}^T y_2) \le 0 + 1 = 1$, so $s(A) = -1 + s(M) \le 0$.
\end{proof}

\begin{remark}
In Case~2 of Lemma~\ref{lem:Petersen}, the singular Petersen interior does not merely preserve the signature bound—it actively absorbs one unit of positive signature through $s(A_{\mathrm{int}}|_S) = -1$, eliminating the need for any sign-controlled edge replacement on $G_2$ and allowing direct induction on the smaller induced subgraph $G_2 - \{x, y\}$.
\end{remark}

% ============================================================================
\section{Proofs of Theorem~\ref{thm:main} and Corollary~\ref{cor:MYL}}\label{sec:proof}
% ============================================================================

We now assemble the reductions of Sections~\ref{sec:prelim}--\ref{sec:schur} via the SPQR-tree decomposition to prove Theorem~\ref{thm:main} and deduce the Ma--Yang--Li conjecture.

\begin{proof}[Proof of Theorem~\ref{thm:main}]
Suppose for contradiction that Theorem~\ref{thm:main} fails, and choose a counterexample $A$ (with underlying NOC signed graph $(G, \sigma)$ and $s(A) > 0$) minimizing $n = |V(G)|$.

If $G$ is disconnected, $A$ is a direct sum of smaller NOC matrices, so $s(A) \le 0$, a contradiction. If $G$ is bipartite, switching on one bipartition class sends $A$ to $-A$, so $s(A) = s(-A) = -s(A)$ and hence $s(A) = 0$, a contradiction. Thus $G$ is connected and non-bipartite, so $n \ge 3$.

By Lemma~\ref{lem:cutvertex} and the minimality of $n$, $G$ has no cut-vertex, so $G$ is $2$-connected. By Lemma~\ref{lem:deg2} and the minimality of $n$, $G$ has no vertex of degree $2$; since $G$ is $2$-connected, $\delta(G) \ge 3$.

If $G$ is $3$-connected, then Theorem~\ref{thm:3conn} implies that $(G, \sigma)$ is switching-equivalent to either $(K_4, +1)$ or $(\mathrm{Petersen}, -1)$. Lemma~\ref{lem:K4}(1) in the first case and Lemma~\ref{lem:Petersen}(1) in the second case yield $s(A) \le 0$, contradicting $s(A) > 0$.

Therefore $G$ is $2$-connected with $\delta(G) \ge 3$, but not $3$-connected. Consider the SPQR tree $\mathcal{T}$ of $G$~\cite{Tut66,DT96}, with $Q$-nodes omitted. Since $G$ is not $3$-connected, $\mathcal{T}$ has at least two nodes and therefore at least two leaves. Examine the possible types of a leaf node $\mu$ of $\mathcal{T}$:
\begin{itemize}
  \item $\mu$ cannot be a $P$-node: a $P$-node consists of two poles joined by at least three parallel edges, and since $G$ is simple at most one of them is a real edge, so at least two are virtual edges, giving $\deg_{\mathcal{T}}(\mu) \ge 2$.
  \item $\mu$ cannot be an $S$-node: a leaf $S$-node is a cycle of length at least $3$ with exactly one virtual edge $xy$, so its remaining edges form a path of length at least $2$ in $G$ whose internal vertices have degree $2$ in $G$, contradicting $\delta(G) \ge 3$.
\end{itemize}
Hence every leaf of $\mathcal{T}$ is an $R$-node, whose skeleton is a $3$-connected simple graph $T = H + xy$ with a single virtual edge $xy$. Thus $H = T - xy$ is a subgraph of $G$ with $xy \notin E(H)$, and $G = H \cup G_2$ with $V(H) \cap V(G_2) = \{x, y\}$ (where $xy$ belongs to $E(G_2)$ if it is a real edge of $G$).

We claim that no leaf $R$-node of $\mathcal{T}$ can have $H$ bipartite. Suppose to the contrary that $H$ is bipartite for some leaf $R$-node $T = H + xy$. By Proposition~\ref{prop:bipartite}, $H$ contains two $(x, y)$-paths $P_1, P_2$ of the same parity with $w_\sigma(P_1) - w_\sigma(P_2) \equiv 2 \pmod 4$. If $G_2$ contained an $(x, y)$-path $Q$ of the opposite parity, then $P_1 \cup Q$ and $P_2 \cup Q$ would be two odd cycles in $G$ whose mod-$4$ weights differ by $2 \pmod 4$, contradicting NOC. Thus every $(x, y)$-path in $G_2$ has the same parity as $P_1, P_2$. Because $G$ is $2$-connected, $G_2 + xy$ is $2$-connected, so $G_2$ is either $2$-connected or a block-chain between $x$ and $y$; in particular, every block of $G_2$ lies on an $(x, y)$-path in $G_2$. If any block $B$ of $G_2$ were non-bipartite, let $s, t \in V(B)$ be the distinct vertices where an $(x, y)$-path enters and leaves $B$, and let $C_{\mathrm{odd}}$ be an odd cycle in $B$. By Menger's theorem applied to the $2$-connected graph $B$, there exist two vertex-disjoint paths in $B$ from $\{s, t\}$ to $C_{\mathrm{odd}}$; combining them with the two arcs of $C_{\mathrm{odd}}$ produces both an even and an odd $(s, t)$-path in $B$, and hence $(x, y)$-paths of both parities in $G_2$, a contradiction. Therefore every block of $G_2$ is bipartite, so $G_2$ is bipartite with $(x, y)$-paths of the same parity as $H$. It follows that $G = H \cup G_2$ is bipartite, contradicting that $G$ is non-bipartite.

Now pick any leaf $R$-node $\mu$ of $\mathcal{T}$, with skeleton $T = H + xy$ and complement $G_2$ (so $G = H \cup G_2$, $V(H) \cap V(G_2) = \{x, y\}$, and $xy \notin E(H)$). By the claim above, $H$ is non-bipartite. Because $\mathcal{T}$ has at least two leaves, there is a second leaf $R$-node $\mu' \ne \mu$ whose subgraph $H'$ lies entirely inside $G_2$. Applying the same claim to $\mu'$ shows that $H'$ is non-bipartite, so $G_2$ is non-bipartite. As shown in the preceding paragraph, a non-bipartite block of $G_2$ yields $(x, y)$-paths of both parities in $G_2$; let $Q$ be an odd $(x, y)$-path in $G_2$.

Define a signing $\sigma'$ on the $3$-connected non-bipartite graph $T = H + xy$ by setting $\sigma'|_H = \sigma|_H$ and choosing $\sigma'(xy) \in \{+1, -1\}$ so that $w_{\sigma'}(xy) \equiv w_\sigma(Q) \pmod 4$. Every odd cycle of $T$ avoiding $xy$ is an odd cycle of $H \subseteq G$, and every odd cycle of $T$ using $xy$ corresponds to an odd cycle of $G$ of the same mod-$4$ weight obtained by replacing $xy$ with $Q$. Since $(G, \sigma)$ is NOC, $(T, \sigma')$ is a $3$-connected non-bipartite NOC signed graph. By Theorem~\ref{thm:3conn}, $(T, \sigma')$ is switching-equivalent to either $(K_4, +1)$ or $(\mathrm{Petersen}, -1)$. Restricting to $H = T - xy$, $(H, \sigma|_H)$ is switching-equivalent to either $(K_4 - xy, +1)$ or $(\mathrm{Petersen} - xy, -1)$. By Lemma~\ref{lem:K4}(2) in the first case and Lemma~\ref{lem:Petersen}(2) in the second case, the minimality of $n$ gives $s(A) \le 0$, contradicting $s(A) > 0$.
\end{proof}

\begin{proof}[Proof of Corollary~\ref{cor:MYL}]
We prove $s(G) \le c_5(G)$ by induction on $c_5(G)$. If $c_5(G) = 0$, take $A = A(G)$ (so $\sigma \equiv +1$). For any odd cycle $C$ of length $2k + 1$, $w_\sigma(C) \equiv 2k + 1 \pmod 4$; since $c_5(G) = 0$, no cycle has length $1 \pmod 4$, so $(G, +1)$ is NOC and Theorem~\ref{thm:main} yields $s(G) \le 0 = c_5(G)$. If $c_5(G) \ge 1$, choose a vertex $v$ belonging to at least one $(1 \bmod 4)$-cycle of $G$. Then $c_5(G - v) \le c_5(G) - 1$. By Cauchy's interlacing theorem, deleting one row and column from a real symmetric matrix satisfies $n^+(G) - 1 \le n^+(G - v) \le n^+(G)$ and $n^-(G) - 1 \le n^-(G - v) \le n^-(G)$, whence $|s(G) - s(G - v)| \le 1$. By the induction hypothesis,
\[
  s(G) \le s(G - v) + 1 \le c_5(G - v) + 1 \le c_5(G).
\]

For the lower bound $-c_3(G) \le s(G)$, we induct on $c_3(G)$. If $c_3(G) = 0$, take $A = -A(G)$ (so $\sigma \equiv -1$). For an odd cycle $C$ of length $2k + 1$, $\epsilon(C) = (-1)^k (-1)^{2k+1} = (-1)^{k+1}$, which equals $+1$ if and only if $2k + 1 \equiv 3 \pmod 4$. Since $c_3(G) = 0$, $(G, -1)$ is NOC, so Theorem~\ref{thm:main} gives $s(-A(G)) \le 0$, i.e., $s(G) = -s(-A(G)) \ge 0$. If $c_3(G) \ge 1$, pick a vertex $v$ on a $(3 \bmod 4)$-cycle of $G$; then $c_3(G - v) \le c_3(G) - 1$ and $s(G) \ge s(G - v) - 1 \ge -c_3(G - v) - 1 \ge -c_3(G)$.
\end{proof}

% ============================================================================
\section*{Use of Artificial Intelligence}
% ============================================================================

The artificial intelligence system Gemini~4 was used during the preparation of this manuscript to assist with proof exploration, symbolic and numerical verification of the matrix identities and Schur-complement calculations, literature review, and LaTeX exposition. All mathematical statements, derivations, and proofs have been reviewed and verified by the author, who bears sole responsibility for the accuracy and completeness of the paper.

% ============================================================================
\begin{thebibliography}{AEKPT26}

\bibitem[AEKPT26]{AEKPT26}
S.~Akbari, C.~Elphick, H.~Kumar, S.~Pragada, and Q.~Tang,
\emph{A new conjecture on the inertia of graphs},
Discrete Math. \textbf{349} (2026), no.~4, 114953;
\url{https://arxiv.org/pdf/2508.01163}.

\bibitem[BGLP26]{BGLP26}
Y.~Bai, A.~Grzesik, B.~Li, and M.~Prorok,
\emph{Cycle lengths in graphs of given minimum degree},
arXiv:2511.03085v2 (2026);
\url{https://arxiv.org/html/2511.03085v2}.

\bibitem[DT96]{DT96}
G.~Di~Battista and R.~Tamassia,
\emph{On-line planarity testing},
SIAM J. Comput. \textbf{25} (1996), no.~5, 956--997.

\bibitem[GLMX24]{GLMX24}
J.~Gao, B.~Li, J.~Ma, and T.~Xie,
\emph{On two cycles of consecutive even lengths},
J. Graph Theory \textbf{106} (2024), no.~2, 225--238;
\url{https://arxiv.org/html/2210.03959v1}.

\bibitem[HJ13]{HJ13}
R.~A.~Horn and C.~R.~Johnson,
\emph{Matrix Analysis}, 2nd ed.,
Cambridge University Press, Cambridge, 2013.

\bibitem[HQSW21]{HQSW21}
C.~He, S.~Qian, H.~Shan, and B.~Wu,
\emph{The signature of chordal graphs and cographs},
Graphs Combin. \textbf{37} (2021), 643--650.

\bibitem[LWZ25]{LWZ25}
H.~Lin, G.~Wang, and W.~Zhou,
\emph{A strengthening on consecutive odd cycles in graphs of given minimum degree},
J. Graph Theory \textbf{110} (2025), no.~4, 431--436;
\url{https://arxiv.org/html/2410.00648v2}.

\bibitem[MYL13]{MYL13}
H.~Ma, W.~Yang, and S.~Li,
\emph{Positive and negative inertia index of a graph},
Linear Algebra Appl. \textbf{438} (2013), no.~1, 331--341.

\bibitem[Tut66]{Tut66}
W.~T.~Tutte,
\emph{Connectivity in Graphs},
Mathematical Expositions, No.~15, University of Toronto Press, Toronto, 1966.

\bibitem[WF14]{WF14}
L.~Wang and Y.-Z.~Fan,
\emph{The signature of line graphs and power trees},
Linear Algebra Appl. \textbf{448} (2014), 264--273;
\url{https://arxiv.org/abs/1310.1003}.

\end{thebibliography}

\vspace{1.2em}
\noindent\textsc{Prince Anand}\\
\textit{E-mail address:} [\texttt{princeanand123pc@gmail.com}]

\end{document}
